\documentclass[11pt,reqno]{amsart}
\usepackage[margin=1.1in]{geometry}
\usepackage{amsmath,amssymb,amsthm}
\usepackage{booktabs}
\usepackage{array}
\usepackage{tikz}
\usepackage[expansion=false]{microtype}
\sloppy
\usepackage[colorlinks=true,linkcolor=blue!50!black,citecolor=blue!50!black,urlcolor=blue!50!black]{hyperref}
\usepackage{aliascnt}
\usepackage[capitalize,nameinlink]{cleveref}
\usepackage{xcolor}
\hypersetup{pdftitle={The Ray--West parameter and the Wilf classification of permutations of length eight},pdfauthor={Mingchang Liu}}

% ---------- theorem environments ----------
\theoremstyle{plain}
\newtheorem{theorem}{Theorem}[section]
\newaliascnt{lemma}{theorem}
\newtheorem{lemma}[lemma]{Lemma}
\aliascntresetthe{lemma}
\newaliascnt{proposition}{theorem}
\newtheorem{proposition}[proposition]{Proposition}
\aliascntresetthe{proposition}
\newaliascnt{corollary}{theorem}
\newtheorem{corollary}[corollary]{Corollary}
\aliascntresetthe{corollary}
\newaliascnt{conjecture}{theorem}
\newtheorem{conjecture}[conjecture]{Conjecture}
\aliascntresetthe{conjecture}
\newaliascnt{question}{theorem}
\newtheorem{question}[question]{Question}
\aliascntresetthe{question}
\theoremstyle{definition}
\newaliascnt{definition}{theorem}
\newtheorem{definition}[definition]{Definition}
\aliascntresetthe{definition}
\newaliascnt{example}{theorem}
\newtheorem{example}[example]{Example}
\aliascntresetthe{example}
\theoremstyle{remark}
\newaliascnt{remark}{theorem}
\newtheorem{remark}[remark]{Remark}
\aliascntresetthe{remark}

% Distinct alias counters preserve one numerical sequence while recording
% the semantic environment type for cleveref and hyperref.
\crefname{lemma}{Lemma}{Lemmas}
\crefname{proposition}{Proposition}{Propositions}
\crefname{corollary}{Corollary}{Corollaries}
\crefname{definition}{Definition}{Definitions}
\crefname{remark}{Remark}{Remarks}
\crefname{appendix}{Appendix}{Appendices}
\Crefname{lemma}{Lemma}{Lemmas}
\Crefname{proposition}{Proposition}{Propositions}
\Crefname{corollary}{Corollary}{Corollaries}
\Crefname{definition}{Definition}{Definitions}
\Crefname{remark}{Remark}{Remarks}
\Crefname{appendix}{Appendix}{Appendices}

% ---------- notation ----------
\newcommand{\Av}{\operatorname{Av}}
\newcommand{\Del}{\operatorname{Del}}
\newcommand{\std}{\operatorname{std}}
\newcommand{\Cat}{\operatorname{Cat}}
\newcommand{\LIS}{\operatorname{LIS}}
\newcommand{\Occ}{\operatorname{Occ}}

\let\greekiota\iota
\renewcommand{\iota}{\boldsymbol{\greekiota}}
\newcommand{\dd}{\boldsymbol{\delta}}
\newcommand{\bb}[1]{\left[#1\right]}
\newcommand{\Sym}{\mathfrak{S}}
\newcommand{\ZZ}{\mathbb{Z}}
\newcommand{\NN}{\mathbb{N}}
\newcommand{\PP}{\mathbb{P}}
\newcommand{\EE}{\mathbb{E}}
\DeclareMathOperator{\Tr}{Tr}
\newcommand{\Gam}{\Gamma}
\newcommand{\bz}{\beta_0}
\newcommand{\bone}{\beta_1}

\begin{document}

\title[The Ray--West parameter and the Wilf classification of $\Sym_8$]{The Ray--West parameter\\ and the Wilf classification\\ of permutations of length eight}

\author{Mingchang Liu}
\email{mliu416@gatech.edu}
\date{}

\subjclass[2020]{05A05, 05A15, 05A19, 05E10}
\keywords{permutation pattern, Wilf equivalence, pattern containment, cluster expansion, Plancherel measure, discrete Bessel kernel}

\begin{abstract}
For a permutation $\tau$ of length $a$, let $g_k(\tau)$ count the
permutations of length $a+k$ containing $\tau$. We express the Ray--West
parameter $j(\tau)=C_2(a)-g_2(\tau)$, where $C_2$ is their universal
codimension-two polynomial, as a statistic of the intervals of $\tau$,
by selecting the primitive supports in Ray and West's insertion
classification. A deletion-graph cycle basis gives a separate
interpretation through one-move classes of occurrences. The interval
formula characterizes the maximizers of $j$ and gives a Poisson limit
law with mean $2$.
We prove that permutations of length eight have exactly $4{,}755$
Wilf classes, generated by the known symmetries and the equivalences of
Backelin--West--Xin and Stankova--West. Exact avoidance counts separate
them through length $14$, and this cutoff is sharp.
For $g_3$ we give a finite formula for arbitrary patterns and a simple
correction for layered patterns. These results also characterize the
entire Wilf class of the increasing pattern in every length.
An elementary cluster expansion organizes the enumeration; its
specialization to increasing patterns yields exact binomial formulas
and Catalan-square corrections.
\end{abstract}

\maketitle

\section{Introduction}\label{sec:intro}

\subsection{Wilf equivalence in low codimension}

Let $\Sym_n$ be the set of permutations of $[n]=\{1,\dots,n\}$, written in one-line notation. A permutation $\pi\in\Sym_n$ \emph{contains} $\tau\in\Sym_a$ if some subsequence of $\pi$ is order-isomorphic to $\tau$; the set of positions of such a subsequence is an \emph{occurrence} of $\tau$ in $\pi$, and $\pi$ \emph{avoids} $\tau$ if it contains no occurrence. Write $\Av_n(\tau)$ for the set of $\tau$-avoiding permutations of length $n$. Two patterns $\tau,\tau'$ are \emph{Wilf-equivalent}, $\tau\sim\tau'$, if $|\Av_n(\tau)|=|\Av_n(\tau')|$ for all $n$. Following Vatter's survey \cite{Vatter2026}, for $k\ge0$ put
\begin{equation}\label{eq:gk}
g_k(\tau):=\#\{\pi\in\Sym_{a+k}:\ \pi\text{ contains }\tau\},\qquad |\Av_n(\tau)|=n!-g_{n-a}(\tau)\quad(n\ge a).
\end{equation}
For $n<a$, every permutation avoids $\tau$, so $|\Av_n(\tau)|=n!$.

The sequence $(g_k(\tau))_{k\ge0}$ determines the Wilf class of $\tau$, and its first terms are necessary conditions for Wilf equivalence that can be evaluated by hand. Pratt \cite{Pratt1973} observed that $g_1(\tau)=a^2+1$ for every $\tau\in\Sym_a$.\footnote{Ray and West \cite[p.~57]{RayWest2003} attribute the same count to an unpublished result of Bloomberg from 1990. Homberger \cite{Homberger2012} gives the run-based argument used in \cref{sec:prelim}.} Ray and West \cite{RayWest2003} proved that
\begin{equation}\label{eq:RW}
g_2(\tau)=C_2(a)-j(\tau),\qquad C_2(a):=\frac{a^4+2a^3+a^2+4a+4}{2},\qquad 0\le j(\tau)\le a-1,
\end{equation}
for an integer $j(\tau)$ depending on $\tau$. 

The survey \cite{Vatter2026} records the numbers $1,1,1,3,16,91,595$ of Wilf classes of $\Sym_m$ for $m\le7$ (sequence A099952 of \cite{OEIS}; Stankova \cite{Stankova1994,Stankova1996} established the classification of length $4$, and Stankova and West \cite{StankovaWest2002} extended it to length $7$) and poses, among others, the following three problems, which we restate; the survey records that the second and third were raised by Vatter at the Permutation Patterns conference of 2007 (see also the problem-session account \cite{Vatter2010}).
\begin{itemize}
\item \textbf{Question 3.1.} How many Wilf-equivalence classes of permutations of length $8$ are there?
\item \textbf{Problem 3.7.} Express the quantity $j$ in the Ray--West formula \eqref{eq:RW} in terms of statistics of $\tau$.
\item \textbf{Problem 3.8.} Find a formula for $g_3(\tau)$.
\end{itemize}

We answer Question 3.1 by a complete finite classification, and Problem 3.7
by an interval statistic derived from Ray and West's codimension-two
classification (\cref{thm:E-intro,thm:F-intro}). For Problem 3.8 we obtain
an explicit finite formula, rather than a comparably compact statistic
of $\tau$ (\cref{thm:G-intro}); the latter problem remains open.

The organizing device is an exact expansion of $g_k(\tau)$ in terms of \emph{clusters}, permutations every point of which lies in an occurrence of $\tau$; it is Möbius inversion over the Boolean lattice of position sets, and its value lies in what it organizes. The expansion isolates, in each codimension $k$, a universal part depending only on $a$ and a correction depending on $\tau$; in codimension two the correction is $j(\tau)$, and the main structural theorem identifies it as a count of \emph{one-move classes} of two-point insertions.

\subsection{Main results}\label{sec:main}

Throughout, $a\ge1$ and $\tau\in\Sym_a$ is fixed. For $\sigma\in\Sym_r$ and $T\subseteq[r]$ let $\sigma|_T$ denote the standardization of the subsequence of $\sigma$ on the positions in $T$. The \emph{cluster weight} of $\sigma$ is
\begin{equation}\label{eq:cw}
c(\sigma):=\sum_{T\subseteq[r]}(-1)^{r-|T|}\,\bb{\sigma|_T\text{ contains }\tau},
\end{equation}
and the \emph{cluster numbers} of $\tau$ are $c_r(\tau):=\sum_{\sigma\in\Sym_r}c(\sigma)$.

\begin{theorem}[cluster expansion; Möbius inversion]\label{thm:A-intro}
For every $\tau\in\Sym_a$ and $k\ge0$,
\begin{equation}\label{eq:clusterexp}
g_k(\tau)=\sum_{r=a}^{a+k}\binom{a+k}{r}^2(a+k-r)!\;c_r(\tau),
\end{equation}
and $c(\sigma)=0$ unless every point of $\sigma$ lies in an occurrence of $\tau$. Moreover $c_a(\tau)=1$, $c_{a+1}(\tau)=-2a$ and $c_{a+2}(\tau)=2a(a+2)-j(\tau)$.
\end{theorem}

The next theorem is the heart of the paper. For $\rho\in\Sym_{a+2}$ the \emph{deletion graph} $\Gam_\tau(\rho)$ has vertex set $[a+2]$ and an edge $\{x,y\}$ whenever deleting the points at positions $x$ and $y$ from $\rho$ leaves a permutation order-isomorphic to $\tau$. Its edges are the occurrences of $\tau$ in $\rho$, and two occurrences are joined by a path if and only if they are connected by a chain of \emph{one-moves}, each replacing a single point of the occurrence. Let $\bz(\Gam_\tau(\rho))$ be the number of connected components of the graph induced on the non-isolated vertices.

\begin{theorem}[one-move classes in codimension two]\label{thm:B-intro}
For every $\tau\in\Sym_a$,
\[
\sum_{\rho\in\Sym_{a+2}}\bz(\Gam_\tau(\rho))=C_2(a).
\]
Consequently the Ray--West parameter is
\begin{equation}\label{eq:jB}
j(\tau)=\sum_{\rho\in\Sym_{a+2},\ \rho\supseteq\tau}\bigl(\bz(\Gam_\tau(\rho))-1\bigr)\;\ge\;0:
\end{equation}
$j(\tau)$ is the number of permutations of length $a+2$ in which $\tau$ occurs in two essentially different ways, each such permutation counted once for every one-move class of occurrences beyond the first.
\end{theorem}

The proof (\cref{sec:delgraph,sec:thmA}) shows that the cycle space of every deletion graph has an explicit basis consisting of triangles and quadrilaterals indexed by inflations of $\tau$, so that the first Betti number of $\Gam_\tau(\rho)$ is an inflation count whose sum over $\rho$ is universal. The same mechanism proves Pratt's theorem: one-point insertions form paths with $(a+1)^2$ vertices and $2a$ edges.

The Ray--West parameter also has an explicit evaluation on $\tau$ itself, which answers Problem 3.7 in the form asked. Primitive returns are the patterns $w\in\Sym_l$, $l\ge2$, that pass the canonical two-point return test of \cref{def:return}, but none of whose proper prefixes with value set $\{1,\dots,t\}$ passes that test.

\begin{theorem}[interval formula]\label{thm:F-intro}
For every $\tau\in\Sym_a$, $j(\tau)$ is the number of pairs $(I,\epsilon)$ where $I$ is an interval of $\tau$ (a set of consecutive positions carrying consecutive values) of size at least two and $\epsilon$ is an orientation, such that the pattern of $\tau$ on $I$, read in the direction $\epsilon$, is a primitive return. The primitive returns of length $l$ are selected from the four-parameter Ray--West track family (\cref{thm:catalogue}) with $r_l=(l^3-l)/12$ members for odd $l$ and $(l^3+2l)/12-[l\ge4]$ for even $l$; every interval of size two counts exactly once. Consequently $0\le j(\tau)\le a-1$, with $j(\tau)=a-1$ exactly for layered and reverse-layered patterns ($2^a-2$ patterns for $a\ge2$), and $j$ of a uniformly random permutation of length $n$ converges in distribution to a Poisson law with mean $2$.
\end{theorem}

\Cref{thm:F-intro} is derived in \cref{sec:interval} from the codimension-two classification of Ray and West, which identifies $j$ as a count of synonymous insertions of a specific type; the interval formula translates that count into a test on the intervals of $\tau$, and the primitive-support count, the substitution-tree recursion behind the extremal statement and the Poisson law follow from the formula.

In codimension three we do not have a statistic as compact as $j$; what we have is an exact finite formula for every pattern, in the form of a count of canonical representatives, and closed forms for special families.

\begin{theorem}[codimension three]\label{thm:G-intro}
For every $\tau\in\Sym_a$, $g_3(\tau)$ is given by the deletion-overlap formula of \cref{thm:boxes}: $6\binom{a+3}3^2$ minus an explicit sum, over the deletion minors of $\tau$ and the relative orders of three inserted points, of sizes of unions of boxes in $\ZZ^3$, evaluated in $O(a^6)$ operations. For layered $\tau$,
\[
g_3(\tau)=g_3(\iota_a)-\#\{\text{interior layers of size}\ge2\}
\]
(\cref{thm:layered}), and $g_3(\tau)=C_3(a)$ for every $\tau$ in which any two points are at Chebyshev distance at least $4$ (\cref{thm:separated}), where $C_3(a)$ is the universal constant of \cref{thm:C-intro} below.
\end{theorem}


The layered correction also determines the full Wilf class of the
increasing pattern in every length $m\ge2$: it has exactly $m(m-1)$
members, and membership is recognized by $g_2$ and $g_3$
(\cref{thm:monoclass}).

The finite classification is the second principal application.

\begin{theorem}[Wilf classes of $\Sym_8$]\label{thm:E-intro}
$\Sym_8$ has exactly $4{,}755$ Wilf classes. Every Wilf equivalence between permutations of length $8$ is generated by the eight symmetries of the square, the equivalences $\iota_k\oplus\gamma\sim\dd_k\oplus\gamma$ of Backelin, West and Xin, and the equivalences $231\oplus\gamma\sim312\oplus\gamma$ of Stankova and West. Two permutations of length $8$ with $|\Av_n(\tau)|=|\Av_n(\tau')|$ for all $n\le14$ are Wilf-equivalent, and $14$ is the least such bound: the numbers of classes distinguished by $|\Av_t|$, $t\le n$, are $1,8,256,4210,4751,4755$ for $n=9,\dots,14$.
\end{theorem}

The upper bound is an exact finite computation with the cited theorems; the lower bound is an exact computation of $|\Av_n(\tau)|$, $n\le14$, for one representative of each class by the cluster expansion. \Cref{sec:S8} describes the reduction, the certificate and the cross-checks, including complete agreement with an earlier computation by two other algorithms, and with a third, brute-force count for all representatives at $n=11,12$ and for the eight representatives whose classes are separated only at $n=14$ at $n=13,14$.

The cluster expansion also isolates a universal part of $g_k$ in every codimension. In higher codimension the natural object is an \emph{inflation complex} $K_\tau(\rho)$ whose vertices are the occurrences of $\tau$ in $\rho\in\Sym_{a+k}$ and whose Euler characteristic $\chi_\tau(\rho)$ sums to a universal constant.

\begin{theorem}[universal constants]\label{thm:C-intro}
Let $N_d(a)$ be the number of ways to inflate the $a$ points of $\tau$ by permutations with $d$ extra points in total, that is, $N_d(a)=[t^d]\bigl(\sum_{i\ge0}(i+1)!\,t^i\bigr)^a$, and put
\begin{equation}\label{eq:Ck}
C_k(a):=\sum_{d=0}^{k}(-1)^d N_d(a)\binom{a+k}{k-d}^2(k-d)!.
\end{equation}
Then $\sum_{\rho\in\Sym_{a+k}}\chi_\tau(\rho)=C_k(a)$ for every $\tau\in\Sym_a$, and
\begin{equation}\label{eq:gkC}
g_k(\tau)=C_k(a)-\sum_{\rho\in\Sym_{a+k},\ \rho\supseteq\tau}\bigl(\chi_\tau(\rho)-1\bigr).
\end{equation}
Explicitly $C_1=a^2+1$, $C_2$ is the Ray--West constant, and
\begin{align*}
C_3(a)&=\frac{a^6+6a^5+10a^4+10a^3+37a^2+44a+36}{6},\\
C_4(a)&=\frac{a^8+12a^7+50a^6+88a^5+137a^4+444a^3+724a^2+848a+576}{24}.
\end{align*}
\end{theorem}

The two leading coefficients agree with the general asymptotic formula
of Ray and West \cite[Cor.~9.5, p.~86]{RayWest2003}, for each fixed $k\ge1$:
\[
g_k(\tau)=\frac{a^{2k}+k(k-1)a^{2k-1}}{k!}
             +O_k(a^{2k-2})\qquad(a\longrightarrow\infty).
\]
Thus the novelty of \cref{thm:C-intro} is the exact all-codimension
Euler-characteristic identity, not these leading coefficients.
For $k=3$ our special-family formulas give a cubic difference between
two families of patterns of the same length (\cref{cor:cubic-spread}).
This contradicts Ray and West's conjectural restriction that the
pattern-dependent terms have degree less than $k$
\cite[p.~57]{RayWest2003}; it does not affect their leading-term estimate.

For $k\le2$ every component of $K_\tau(\rho)$ has Euler characteristic one (\cref{thm:pratt,prop:chi2}), so $\chi_\tau(\rho)$ is the number of one-move classes and \eqref{eq:gkC} is \cref{thm:B-intro}. An exhaustive computation of the Euler characteristic of every component shows that the same holds for $k=3$ and all $\tau$ of length $\le6$, and for $k=4$ and all $\tau$ of length $\le5$, so that in these ranges $g_3(\tau)=C_3(a)-\Delta_3(\tau)$ with $\Delta_3(\tau)$ the number of extra one-move classes of three-point insertions (\cref{thm:k34}); we conjecture this for all $a$, but it fails for $k=5$ (\cref{prop:k5}). The one-move description therefore has a definite boundary, and the exact identity $g_3=C_3(a)-(C_3(a)-g_3)$ should not be read as a proved interpretation of the correction beyond these ranges.

Finally, for the increasing pattern $\iota_a=12\cdots a$ the cluster numbers can be computed exactly.

\begin{theorem}[cluster numbers of the increasing pattern]\label{thm:D-intro}
Let $P_\theta$ be the Poissonized Plancherel measure on partitions, $P_\theta(\lambda)=e^{-\theta}\theta^{|\lambda|}(f^\lambda/|\lambda|!)^2$. Then
\begin{equation}\label{eq:plancherel-intro}
\sum_{r\ge a}c_r(\iota_a)\,\frac{\theta^r}{r!^2}=P_\theta(\lambda_1\ge a),
\end{equation}
and the Fredholm expansion of the discrete Bessel kernel gives, with $\Cat(n)=\binom{2n}{n}/(n+1)$,
\[
c_{a+k}(\iota_a)=(-1)^k\binom{2a+2k-2}{k}\quad(0\le k\le a+1),\qquad
c_{2a+2}(\iota_a)=(-1)^a\binom{4a+2}{a+2}-\Cat(a+1)^2 .
\]
In particular, for $a\le n\le 2a+1$ the number of permutations of $[n]$ with an increasing subsequence of length $a$ equals $\sum_{r=a}^{n}(-1)^{r-a}\binom nr^2(n-r)!\binom{2r-2}{r-a}$.
\end{theorem}

\subsection{Dependency boundary}\label{sec:boundary}

The proofs use the following published results: the run characterization
of deletions, in the form given by Homberger \cite[preprint, Lem.~5]{Homberger2012},
and Pratt's one-point extension theorem \cite{Pratt1973} (both reproved in
\cref{sec:prelim}); the codimension-two classification of Ray and West
\cite{RayWest2003} (\cref{thm:RW}, used to derive the interval formula);
the Wilf equivalences of Backelin--West--Xin \cite{BWX2007} and
Stankova--West \cite{StankovaWest2002}; Schensted's theorem
\cite{Schensted1961} and the hook-length formula \cite{FRT1954};
the determinantal structure of Poissonized Plancherel measure
\cite{BOO2000,Johansson2001}; and substitution decomposition
\cite{BHV2008}. The classification in \cref{thm:E-intro} includes finite
computations for both its upper and lower bounds, described in \cref{sec:S8}.
\Cref{thm:k34,prop:k5} are also computational, in the explicitly stated
finite ranges. The structural proofs of \cref{thm:B-intro,thm:F-intro}
do not depend on the finite classification.

\subsection{Organization}
\Cref{sec:prelim} proves the run lemma and Pratt's theorem. \Cref{sec:cluster} proves the cluster expansion, its generating-function form and the universal constants (\cref{thm:C-intro}, whose proof is a regrouping of the expansion). \Cref{sec:delgraph} studies deletion graphs and proves the cycle-basis lemma; \cref{sec:thmA} deduces \cref{thm:B-intro}. \Cref{sec:interval} proves the interval formula and its consequences, and \cref{sec:codim3} the evaluations of $g_3$ (\cref{thm:G-intro}), with the proof of the layered formula in \cref{app:layered}. \Cref{sec:S8} proves \cref{thm:E-intro} and records the cutoffs for lengths $4$ to $8$. The two remaining sections develop the expansion further: \cref{sec:higher} treats the inflation complex, with the host-by-host form of the corrections, the separated patterns, the computational results for $k=3,4$ and the failure at $k=5$, and \cref{sec:monotone} proves \cref{thm:D-intro}. \Cref{sec:questions} lists questions.

\subsection*{Conventions}
A permutation $\pi\in\Sym_n$ is identified with its set of \emph{points} $(i,\pi(i))$; we refer to a point by its position when no confusion can arise. For a set $D$ of positions, $\pi\setminus D:=\pi|_{[n]\setminus D}$. A \emph{bond} is a pair of adjacent positions carrying adjacent values. A \emph{run} (monotone interval) is a set of consecutive positions carrying consecutive values in increasing or decreasing order; a single point is a run; a \emph{maximal run} is one not contained in a larger run. Every point lies in exactly one maximal run: two runs of size at least two through a point $x$ in different directions would give an ascent bond and a descent bond at $x$, which either use the same neighbour of $x$, which cannot carry two values, or use the two neighbours of $x$, which would then carry the same value; and two runs through $x$ in the same direction have a run as their union, since in both the value of a point is determined by its offset from $x$. 

An \emph{interval} of $\pi$ is a set of consecutive positions carrying consecutive values, in any order. The \emph{direct sum} $\alpha\oplus\beta$ places $\beta$ to the right of and above $\alpha$; the \emph{skew sum} $\alpha\ominus\beta$ places $\beta$ to the right of and below $\alpha$; $\iota_k$ and $\dd_k$ denote the increasing and decreasing permutations of length $k$. The \emph{inflation} $\pi[\sigma_1,\dots,\sigma_n]$ replaces the $i$th point of $\pi$ by a copy of $\sigma_i$ occupying an interval. \par\smallskip\noindent\emph{Symmetries.} 

The eight symmetries are generated by reverse ($r$), complement ($c$) and inverse ($i$). Binomial coefficients $\binom nk$ vanish for $k<0$ or $k>n$.

\section{One-point extensions}\label{sec:prelim}

Fix $\tau\in\Sym_a$. For $\sigma\in\Sym_{a+1}$ containing $\tau$ put $\Del(\sigma):=\{y\in[a+1]:\sigma\setminus y=\tau\}$.
The following run characterization and the resulting proof of the
one-point extension count are given by Homberger
\cite[preprint, Lem.~5 and Cor.~8]{Homberger2012}; the count itself is due to Pratt
\cite{Pratt1973}. We include the arguments because runs will also control
the edges of the deletion graph.

\begin{lemma}[run lemma]\label{lem:run}
If $\sigma\in\Sym_{a+1}$ contains $\tau$, then $\Del(\sigma)$ is a maximal run of $\sigma$.
\end{lemma}

\begin{proof}
Let $y<y'$ be in $\Del(\sigma)$ and let $m_1<\dots<m_t$ be the positions strictly between them. The permutations $\sigma\setminus y$ and $\sigma\setminus y'$ are both equal to $\tau$. Comparing them position by position, they agree outside the block of positions $\{y,m_1,\dots,m_t,y'\}$, and on this block the sequence $(m_1,\dots,m_t,y')$ in $\sigma\setminus y$ occupies the same slots as $(y,m_1,\dots,m_t)$ in $\sigma\setminus y'$. Equality of the two patterns therefore says two things. First, every point $z$ outside the block compares in the same way with the $i$th and the $(i+1)$st entry of $(y,m_1,\dots,m_t,y')$ for every $i$; hence no outside value separates two values of the block, and the block carries consecutive values. Second, for all $i$, $[\sigma(s_i)<\sigma(s_{i+1})]=[\sigma(s_{i+1})<\sigma(s_{i+2})]$ where $s_1,\dots,s_{t+2}$ enumerate the block; hence the block is monotone. So $y$ and $y'$ lie in a common run. Conversely, deleting any point of a run gives the same permutation, so $\Del(\sigma)$ is a union of maximal runs, and by the first part it is a single one.
\end{proof}

\begin{theorem}[Pratt \cite{Pratt1973}]\label{thm:pratt}
$g_1(\tau)=a^2+1$ for every $\tau\in\Sym_a$.
\end{theorem}

\begin{proof}
Count pairs $(\sigma,y)$ with $y\in\Del(\sigma)$: each arises from inserting one point into $\tau$ at one of $(a+1)^2$ slots (a position gap and a value gap), so $\sum_\sigma|\Del(\sigma)|=(a+1)^2$. Next count pairs $(\sigma,\{y,y'\})$ with $\{y,y'\}$ a bond of $\sigma$ contained in $\Del(\sigma)$. By \cref{lem:run} these are the $|\Del(\sigma)|-1$ consecutive pairs of the run $\Del(\sigma)$; and such a pair arises exactly by replacing one point of $\tau$ by a bond in one of two directions, so $\sum_\sigma(|\Del(\sigma)|-1)=2a$. Subtracting, the number of $\sigma$ is $(a+1)^2-2a$.
\end{proof}

The proof exhibits the set of one-point insertions as a disjoint union of paths (slide the inserted point along a run) with $(a+1)^2$ vertices and $2a$ edges and no cycles. For two inserted points there are cycles; \cref{sec:delgraph} describes them.

\begin{lemma}[twins]\label{lem:twins}
If $x,x'$ lie in a common run of $\rho\in\Sym_n$, then $\rho\setminus x=\rho\setminus x'$, and more generally $\rho\setminus(D\cup\{x\})=\rho\setminus(D\cup\{x'\})$ for every set $D$ of positions not containing $x,x'$.
\end{lemma}

\begin{proof}
Deleting either point of a run leaves the remaining points of the run in the same relative position with respect to every other point, and the points of $D$ are among the other points.
\end{proof}

\section{The cluster expansion}\label{sec:cluster}

For $\pi\in\Sym_n$ and $S\subseteq[n]$ write $f(S):=\bb{\pi|_S\text{ contains }\tau}$. Möbius inversion over the Boolean lattice of position sets gives
\begin{equation}\label{eq:mobius}
f(S)=\sum_{T\subseteq S}c(\pi|_T),
\end{equation}
where $c$ is the cluster weight \eqref{eq:cw}; note that $c(\pi|_T)$ depends only on the pattern $\pi|_T$. A permutation $\sigma$ is a \emph{cluster} of $\tau$ if every point of $\sigma$ lies in an occurrence of $\tau$.

\begin{theorem}[cluster expansion; Möbius inversion]\label{thm:cluster}
For every $\tau\in\Sym_a$ and $k\ge0$, identity \eqref{eq:clusterexp} holds, and $c(\sigma)=0$ unless $\sigma$ is a cluster.
\end{theorem}

\begin{proof}
By \eqref{eq:mobius} with $S=[a+k]$, $g_k(\tau)=\sum_{\pi\in\Sym_{a+k}}\sum_{S\subseteq[a+k]}c(\pi|_S)$. The number of pairs $(\pi,S)$ with $|S|=r$ and $\pi|_S=\sigma$ is $\binom{a+k}{r}^2(a+k-r)!$: choose the positions of $S$, the values of $S$, and the arrangement of the remaining $a+k-r$ points. Grouping by $\sigma$ gives \eqref{eq:clusterexp}. If a point $p$ of $\sigma$ lies in no occurrence of $\tau$, then $T\mapsto T\,\triangle\,\{p\}$ is an involution on subsets of $[r]$ that reverses the sign $(-1)^{r-|T|}$ and preserves $\bb{\sigma|_T\supseteq\tau}$, so $c(\sigma)=0$.
\end{proof}

\begin{proposition}\label{prop:firstclusters}
$c_a(\tau)=1$, $c_{a+1}(\tau)=-2a$ and $c_{a+2}(\tau)=2a(a+2)-j(\tau)$, where $j(\tau):=C_2(a)-g_2(\tau)$ is the Ray--West parameter.
\end{proposition}

\begin{proof}
For $r=a$ only $\sigma=\tau$ and $T=[a]$ contribute. For $r=a+1$, $c(\sigma)=\bb{\sigma\supseteq\tau}-|\Del(\sigma)|$, and summing over $\sigma$ gives $g_1(\tau)-(a+1)^2=-2a$ by \cref{thm:pratt}. For $r=a+2$, \eqref{eq:clusterexp} with $k=2$ reads $g_2=2\binom{a+2}{2}^2-2a(a+2)^2+c_{a+2}$, and $2\binom{a+2}2^2-2a(a+2)^2=C_2(a)-2a(a+2)$ by direct expansion.
\end{proof}

Thus the Ray--West parameter is the deviation of the third cluster number from $2a(a+2)$. \Cref{sec:delgraph,sec:thmA} explain this value: $2a(a+2)=6a+4\binom a2$ counts the \emph{inflation structures} of size $a+2$, in which one point of $\tau$ is replaced by a pattern of length three ($6a$ ways) or two points are replaced by bonds ($4\binom a2$ ways), and \cref{prop:chi2} shows that $c(\sigma)$ falls short of the number of inflation structures of a host $\sigma\in\Sym_{a+2}$ by exactly the number of one-move classes of occurrences beyond the first. (The deficit is not confined to hosts that are not inflations of $\tau$; see the examples after \cref{prop:chi2}.) Independently of $\tau$, \eqref{eq:clusterexp} and \cref{prop:firstclusters} give
\begin{equation}\label{eq:g3cluster}
g_3(\tau)=U_3(a)-(a+3)^2\,j(\tau)+c_{a+3}(\tau),\qquad U_3(a):=6\binom{a+3}3^2-4a\binom{a+3}2^2+2a(a+2)(a+3)^2,
\end{equation}
so that Problem 3.8 is the problem of the fourth cluster number.

The expansion inverts: $c_r(\tau)=\sum_{t=a}^{r}(-1)^{r-t}\binom rt^2(r-t)!\,g_{t-a}(\tau)$, and it has a compact generating-function form. Since $\binom nr^2(n-r)!/n!^2=1/\bigl(r!^2(n-r)!\bigr)$, multiplying \eqref{eq:clusterexp} by $\theta^{n}/n!^2$ with $n=a+k$ and summing over $k$ gives
\begin{equation}\label{eq:egf}
\sum_{n\ge a}g_{n-a}(\tau)\frac{\theta^n}{n!^2}=e^{\theta}\sum_{r\ge a}c_r(\tau)\frac{\theta^r}{r!^2},
\qquad\text{i.e.}\qquad
\sum_{r\ge a}c_r(\tau)\frac{\theta^r}{r!^2}=1-e^{-\theta}\sum_{n\ge0}|\Av_n(\tau)|\frac{\theta^n}{n!^2}.
\end{equation}
(The second form uses $|\Av_n(\tau)|=n!$ for $n<a$.) \Cref{sec:monotone} exploits \eqref{eq:egf} for $\tau=\iota_a$.

\begin{remark}
The expansion is a positional analogue of the Goulden--Jackson cluster method for consecutive patterns \cite{GouldenJackson1979}: there clusters are words covered by overlapping consecutive occurrences; here they are permutations covered by occurrences in the pattern-containment order, and the weights are Möbius-type sums over position sets rather than over overlap structures.
\end{remark}

\subsection{Universal constants}\label{sec:universal}

The expansion isolates, in every codimension, a part of $c_r(\tau)$ that depends only on $a$. For $\sigma\in\Sym_r$ let $I_\tau(\sigma)$ be the number of ways to write $\sigma$ as an inflation $\tau[\pi_1,\dots,\pi_a]$ (so $I_\tau(\tau)=1$, and $I_\tau(\sigma)=0$ unless $\sigma$ has $\tau$ as a quotient by intervals), and for $\rho\in\Sym_{a+k}$ put
\begin{equation}\label{eq:chi}
\chi_\tau(\rho):=\sum_{T\subseteq[a+k]}(-1)^{|T|-a}\,I_\tau(\rho|_T).
\end{equation}
\Cref{sec:higher} interprets $\chi_\tau(\rho)$ as the Euler characteristic of a cell complex whose vertices are the occurrences of $\tau$ in $\rho$; here we only need the definition.

\begin{theorem}\label{thm:universal}
With $N_d(a)$ and $C_k(a)$ as in \eqref{eq:Ck}: $\sum_{\rho\in\Sym_{a+k}}\chi_\tau(\rho)=C_k(a)$ for every $\tau\in\Sym_a$, and \eqref{eq:gkC} holds.
\end{theorem}

\begin{proof}
Grouping the pairs $(\rho,T)$ by the pattern $\sigma=\rho|_T$ as in the proof of \cref{thm:cluster},
\[
\sum_{\rho\in\Sym_{a+k}}\chi_\tau(\rho)=\sum_{r=a}^{a+k}\binom{a+k}r^2(a+k-r)!\,(-1)^{r-a}\sum_{\sigma\in\Sym_r}I_\tau(\sigma),
\]
and $\sum_{\sigma\in\Sym_r}I_\tau(\sigma)$ is the number of tuples $(\pi_1,\dots,\pi_a)$ of permutations with $\sum|\pi_i|=r$, which is $[t^{r}]\bigl(\sum_{s\ge1}s!\,t^s\bigr)^a=N_{r-a}(a)$. Identity \eqref{eq:gkC} is $g_k(\tau)=\sum_{\rho\supseteq\tau}1$.
\end{proof}

The first values are $N_1=2a$, $N_2=2a(a+2)$, $N_3=24a+12a(a-1)+8\binom a3$ and $N_4=120a+48a(a-1)+36\binom a2+24a\binom{a-1}2+16\binom a4$ (one block receiving all extra points, or two blocks receiving $3+1$ or $2+2$, or three blocks receiving $2+1+1$, or four blocks receiving one each), which give $C_1,\dots,C_4$ as displayed in \cref{thm:C-intro}. Thus $c_{a+d}(\tau)=(-1)^dN_d(a)+D_{a+d}(\tau)$ with a correction $D_{a+d}(\tau)$ that vanishes for $d\le1$ and equals $-j(\tau)$ for $d=2$ (\cref{prop:firstclusters}); \cref{prop:weights} gives the exact host-by-host form of $D_{a+d}$, and it is not simply the contribution of the hosts that are not inflations of $\tau$.

\section{Two-point extensions: the deletion graph}\label{sec:delgraph}

Fix $\tau\in\Sym_a$ and $\rho\in\Sym_{a+2}$. The \emph{deletion graph} $\Gam=\Gam_\tau(\rho)$ has vertex set $[a+2]$ and edge set
\[
E(\Gam):=\bigl\{\{x,y\}:\ \rho\setminus\{x,y\}=\tau\bigr\}.
\]
Its edges are in bijection with the occurrences of $\tau$ in $\rho$: an occurrence is the complement of its deletion pair. Let $V(\Gam)$ be the set of non-isolated vertices, $\bz(\Gam)$ the number of connected components of the graph induced on $V(\Gam)$ (so $\bz=0$ when $\rho$ avoids $\tau$), and $\bone(\Gam):=|E(\Gam)|-|V(\Gam)|+\bz(\Gam)$ its cycle rank, the dimension of its cycle space over $\mathrm{GF}(2)$. For $x\in V(\Gam)$ put
\[
R_x:=N_\Gam(x)=\{y:\rho\setminus\{x,y\}=\tau\}=\Del(\rho\setminus x).
\]
Two occurrences are related by a \emph{one-move} if they differ in exactly one point, i.e. their deletion pairs share a point; the components of the line graph of $\Gam$, equivalently of $\Gam$ itself restricted to $V(\Gam)$, are the \emph{one-move classes} of occurrences. By \cref{lem:run}, a one-move slides a point along a run of $\rho$ minus the common deleted point.

\begin{lemma}\label{lem:Rx}
For $x\in V(\Gam)$, $R_x$ is a maximal run of $\rho\setminus x$. In particular $R_x$ is a set of positions of $\rho$ that is consecutive except possibly for skipping $x$, and two points of $R_x$ that are adjacent in $R_x$ have values differing by $1$ in $\rho$, or differing by $2$ with the value of $x$ between them.
\end{lemma}

\begin{proof}
Apply \cref{lem:run} to $\sigma=\rho\setminus x$, which contains $\tau$ because $x$ is not isolated.
\end{proof}

\begin{lemma}\label{lem:twinsG}
If $x,x'$ lie in a common run of $\rho$, then for every $y\notin\{x,x'\}$, $\{x,y\}\in E(\Gam)$ if and only if $\{x',y\}\in E(\Gam)$.
\end{lemma}

\begin{proof}
This is \cref{lem:twins} with $D=\{y\}$.
\end{proof}

\begin{definition}[inflation cycles]\label{def:cycles}
Let $T(\rho)$ be the set of triples of consecutive positions $\{i,i+1,i+2\}$ carrying three consecutive values (a \emph{$3$-interval}) all three of whose pairs are edges of $\Gam$, and let $S(\rho)$ be the set of pairs of disjoint bonds $\{x,x+1\}$, $\{y,y+1\}$ with $x+1<y$ all four of whose cross pairs $\{x,y\},\{x,y+1\},\{x+1,y\},\{x+1,y+1\}$ are edges. Put $t(\rho):=|T(\rho)|$ and $s(\rho):=|S(\rho)|$. Each element of $T(\rho)$ is a triangle of $\Gam$ and each element of $S(\rho)$ is a $4$-cycle of $\Gam$.
\end{definition}

\begin{figure}[htbp]
\centering
\begin{tikzpicture}[scale=.7, every node/.style={font=\small}]
\begin{scope}
\coordinate (a) at (0,0); \coordinate (b) at (3,0); \coordinate (c) at (1.5,2);
\draw (a)--(b)--(c)--(a); \draw[line width=1.5pt] (b)--(c);
\foreach \v in {a,b,c} \fill (\v) circle (2pt);
\node[below] at (a) {$i$}; \node[below] at (b) {$i+1$};
\node[above] at (c) {$i+2$};
\node at (1.5,-1) {a $3$-interval};
\end{scope}
\begin{scope}[xshift=7cm]
\coordinate (a) at (0,0); \coordinate (b) at (0,2);
\coordinate (c) at (3,0); \coordinate (d) at (3,2);
\draw (a)--(c)--(b)--(d)--(a); \draw[line width=1.5pt] (b)--(d);
\foreach \v in {a,b,c,d} \fill (\v) circle (2pt);
\node[left] at (a) {$x$}; \node[left] at (b) {$x+1$};
\node[right] at (c) {$y$}; \node[right] at (d) {$y+1$};
\node at (1.5,-1) {two disjoint bonds};
\end{scope}
\end{tikzpicture}
\caption{The inflation cycles of \cref{def:cycles}. Vertices denote
positions in the host. The thick edge is the lexicographically largest
edge of the cycle. Only the cycle edges are drawn; crossings are not
vertices, and additional deletion-graph edges may be present.}
\label{fig:cycles}
\end{figure}

The pairs $(\rho,I)$ with $I\in T(\rho)$ correspond bijectively to pairs $(x,\pi)$ with $x$ a point of $\tau$ and $\pi\in\Sym_3$: contracting $I$ to a point recovers $\tau$ with the marked point $x$ (deleting any two points of $I$ leaves $\tau$), and conversely $\rho=\tau[x\leftarrow\pi]$ has $I$ as the inflated block, all of whose pairs are edges. Likewise the pairs $(\rho,\{B,B'\})$ with $\{B,B'\}\in S(\rho)$ correspond to pairs of distinct points of $\tau$ each replaced by a bond in one of two directions. Hence
\begin{equation}\label{eq:ts}
\sum_{\rho\in\Sym_{a+2}}t(\rho)=6a,\qquad\sum_{\rho}s(\rho)=4\binom a2,\qquad \sum_\rho\bigl(t(\rho)+s(\rho)\bigr)=2a(a+2).
\end{equation}

Order the edges of $\Gam$ lexicographically as sorted pairs: $\{u,v\}<\{u',v'\}$ (with $u<v$, $u'<v'$) if $u<u'$, or $u=u'$ and $v<v'$. The \emph{leading edge} of a triangle $\{i,i+1,i+2\}$ is $\{i+1,i+2\}$, and that of a square $\{x,x+1\}\times\{y,y+1\}$ is $\{x+1,y+1\}$. In both cases the leading edge is the largest edge of the cycle, and distinct elements of $T(\rho)\cup S(\rho)$ have distinct leading edges: the leading edge of a triangle joins adjacent positions, that of a square does not, and each determines its cycle.

\begin{lemma}[leading-edge lemma]\label{lem:leading}
Let $\{u,v\}\in E(\Gam)$ with $u<v$, and suppose that $u$ and $v$ are joined by a path in $\Gam$ all of whose edges are smaller than $\{u,v\}$. Then $u\ge2$ and, writing $u^-:=u-1$ and $v^-:=v-1$:
\begin{enumerate}
\item[$(\alpha)$] if $v=u+1$, then $\{u^-,u,v\}$ is a $3$-interval of $\rho$ all of whose pairs are edges, i.e. $\{u,v\}$ is the leading edge of an element of $T(\rho)$;
\item[$(\beta)$] if $v\ge u+2$, then $\{u^-,u\}$ and $\{v^-,v\}$ are bonds of $\rho$ and all four cross pairs are edges, i.e. $\{u,v\}$ is the leading edge of an element of $S(\rho)$.
\end{enumerate}
\end{lemma}

\begin{proof}
Let the path be $u=p_0,p_1,\dots,p_m=v$; it does not use the edge $\{u,v\}$, so $m\ge2$. The last edge $\{p_{m-1},v\}$ is smaller than $\{u,v\}$; if $p_{m-1}>v$ the sorted pair $\{v,p_{m-1}\}$ would be larger, and $p_{m-1}=u$ is excluded, so $w:=p_{m-1}<u$. In particular $u\ge2$. The first edge $\{u,p_1\}$ is smaller than $\{u,v\}$, so $w':=p_1<v$ (either $w'<u$ or $u<w'<v$). Since $w,u\in R_v$ and $R_v$ is, by \cref{lem:Rx}, a set of consecutive positions of $\rho$ apart from possibly skipping $v$, and $w<u<v$, every position in $[w,u]$ lies in $R_v$; thus $u^-\in R_v$, and $u^-,u$ are adjacent in $R_v$. Likewise, since $w',v\in R_u$, every position in $[w',v]$ other than $u$ lies in $R_u$.

$(\alpha)$ Let $v=u+1$. Then $w'\le u-1$, so $u^-\in R_u$, and $u^-,v$ are adjacent in $R_u$ (they are the positions $u-1$ and $u+1$ of $\rho$, adjacent once $u$ is removed). Let $A,B,C$ be the values of $u^-,u,v$. By \cref{lem:Rx} applied to $R_v$ and to $R_u$: $|A-B|=1$, or $|A-B|=2$ with $C$ between $A$ and $B$; and $|A-C|=1$, or $|A-C|=2$ with $B$ between $A$ and $C$. In every case $\{A,B,C\}$ consists of three consecutive integers: if $|A-B|=|A-C|=1$ then $B$ and $C$ lie on opposite sides of $A$; if $|A-B|=2$ with $C$ between, then $B,C$ lie on the same side of $A$ at distances $2$ and $1$; symmetrically if $|A-C|=2$; and both distances cannot be $2$. So $\{u^-,u,v\}$ is a $3$-interval, and its three pairs are edges: $\{u,v\}$ by hypothesis, $\{u^-,v\}$ because $u^-\in R_v$, and $\{u^-,u\}$ because $u^-\in R_u$.

$(\beta)$ Let $v\ge u+2$. Now $v-1\ne u$, so $v^-\in R_u$ and $v^-,v$ are adjacent in $R_u$. Write $c$ for the value of $u$.

\emph{Claim 1: $\{v^-,v\}$ is a bond.} Suppose not. By \cref{lem:Rx} the values of $v^-$ and $v$ are $c-1$ and $c+1$ in some order; write the value of $v$ as $c+\varepsilon$ with $\varepsilon=\pm1$, so $v^-$ has value $c-\varepsilon$. The value of $u^-$ is adjacent to $c$ in $\rho\setminus v$, hence equals $c-\varepsilon$ (impossible: that value belongs to $v^-\ne u^-$, as $v-1\ge u+1$), or $c+\varepsilon$ (impossible: that value belongs to $v$), or $c+2\varepsilon$. So $u^-$ has value $c+2\varepsilon$. In $\rho\setminus u$ the run $R_u$ contains $v^-,v$ as adjacent points whose values, after removing $c$, are consecutive with $v$ on the side of $\varepsilon$; so $R_u$ is monotone with direction $\varepsilon$. If $w'<u$, then $u^-\in R_u$ precedes $v^-$ in $R_u$, and its value $c+2\varepsilon$ would have to lie on the side opposite to $\varepsilon$ of the value $c-\varepsilon$ of $v^-$, which is false. Hence $u<w'<v$. Then $m\ge3$ (as $p_{m-1}=w<u<w'=p_1$), and the second edge $\{w',p_2\}$ of the path is smaller than $\{u,v\}$; since $w'>u$, this forces $x:=p_2<u$. Thus $x,u\in R_{w'}$, so every position in $[x,u]$ lies in $R_{w'}$; in particular $u^-\in R_{w'}$ and $u^-,u$ are adjacent in $R_{w'}$, whence their values $c+2\varepsilon$ and $c$ are adjacent in $\rho\setminus w'$, i.e. the value $c+\varepsilon$ between them belongs to $w'$. But $c+\varepsilon$ is the value of $v\ne w'$, a contradiction.

\emph{Claim 2: $\{u^-,u\}$ is a bond.} Suppose not. Since $u^-,u$ are adjacent in $R_v$, the value of $u^-$ is $c+2\varepsilon'$ with the value $c+\varepsilon'$ of $v$ between them. By Claim 1, $v^-$ has value $c+\varepsilon'\pm1\in\{c,c+2\varepsilon'\}$, i.e. $v^-\in\{u,u^-\}$, contradicting $v-1\ge u+1$.

Finally, $u^-$ and $u$ are twins and $v^-$ and $v$ are twins in the sense of \cref{lem:twinsG}, so from the edge $\{u,v\}$ we obtain the edges $\{u^-,v\}$, $\{u,v^-\}$ and $\{u^-,v^-\}$. Since $u<v-1$, the two bonds are disjoint and $\{u,v\}$ is the leading edge of the square $\{u^-,u\}\times\{v^-,v\}\in S(\rho)$.
\end{proof}

\begin{lemma}[cycle basis]\label{lem:basis}
The edge sets of the triangles in $T(\rho)$ and the squares in $S(\rho)$ form a basis of the cycle space of $\Gam_\tau(\rho)$ over $\mathrm{GF}(2)$. Consequently
\begin{equation}\label{eq:b1}
\bone(\Gam_\tau(\rho))=t(\rho)+s(\rho).
\end{equation}
\end{lemma}

\begin{proof}
\emph{Independence.} In a nontrivial $\mathrm{GF}(2)$-combination of these cycles, consider the cycle $g$ whose leading edge $e$ is largest. Every other cycle in the combination has all its edges at most its own leading edge, which is smaller than $e$; so $e$ occurs in $g$ only, and the sum is nonzero.

\emph{Spanning.} Let $z\ne0$ be an even subgraph of $\Gam$ and $e=\{u,v\}$ its largest edge. An even subgraph is an edge-disjoint union of cycles, so $e$ lies on a cycle of $z$ all of whose other edges are smaller than $e$; that cycle minus $e$ is a path from $u$ to $v$ as in \cref{lem:leading}, which provides $g\in T(\rho)\cup S(\rho)$ with leading edge $e$. Then $z+\partial g$ is an even subgraph whose largest edge is smaller than $e$. Induction on the largest edge finishes the proof.
\end{proof}

\section{Proof of \texorpdfstring{\cref{thm:B-intro}}{Theorem 1.2}}\label{sec:thmA}

\begin{theorem}\label{thm:A}
For every $\tau\in\Sym_a$, $\sum_{\rho\in\Sym_{a+2}}\bz(\Gam_\tau(\rho))=C_2(a)$.
\end{theorem}

\begin{proof}
For each $\rho$, $\bz=|V|-|E|+\bone$. Summing over $\rho$: $\sum_\rho|E(\Gam)|$ is the number of pairs $(\rho,\{x,y\})$ with $\rho\setminus\{x,y\}=\tau$, i.e. the number of ways to insert two points into $\tau$, which is $2\binom{a+2}2^2$ (choose two position gaps, two value gaps, and the arrangement of the two new values). $\sum_\rho|V(\Gam)|$ is the number of pairs $(\rho,x)$ with $\rho\setminus x\supseteq\tau$, i.e. $(a+2)^2$ times the number of $\sigma\in\Sym_{a+1}$ containing $\tau$, which is $(a+2)^2(a^2+1)$ by \cref{thm:pratt}. And $\sum_\rho\bone=2a(a+2)$ by \eqref{eq:b1} and \eqref{eq:ts}. Hence
\[
\sum_\rho\bz=(a+2)^2(a^2+1)-\frac{(a+2)^2(a+1)^2}{2}+2a(a+2)=\frac{(a+2)^2(a-1)^2}{2}+2a(a+2),
\]
and $(a+2)^2(a-1)^2+4a(a+2)=a^4+2a^3+a^2+4a+4$.
\end{proof}

\begin{corollary}\label{cor:j}
For every $\tau\in\Sym_a$, \eqref{eq:jB} holds; in particular $j(\tau)\ge0$, with $j(\tau)=0$ if and only if in every permutation of length $a+2$ all occurrences of $\tau$ are connected by one-moves.
\end{corollary}

\begin{proof}
$g_2(\tau)$ is the number of $\rho\in\Sym_{a+2}$ with $\bz(\Gam_\tau(\rho))\ge1$. Subtract from \cref{thm:A}.
\end{proof}

\Cref{cor:j} recovers the lower bound in \eqref{eq:RW} and gives an intrinsic meaning to $j$. The upper bound $j\le a-1$ follows from the interval formula of \cref{sec:interval} (or from \cite{RayWest2003}).

\begin{remark}[canonical insertions]\label{rem:canonical}
A second route to \cref{thm:A} would count \emph{canonical} edges. Call an edge $\{z_1<z_2\}$ locally least if there is no edge $\{z_1,y\}$ with $y<z_2$ and no edge $\{z_2,y\}$ with $y<z_1$; by \cref{lem:Rx} this means that $z_2$ is the leftmost point of its run in $\rho\setminus z_1$ and $z_1$ the leftmost point of its run in $\rho\setminus z_2$. The number of locally least edges, summed over $\rho$, is $C_2(a)+a-1$ for every $\tau$. Indeed, an edge is a pair of points inserted into $\tau$, at slots $(i_1,v_1)$ and $(i_2,v_2)$ (a position gap and a value gap, $1\le i,v\le a+1$), and $z_2$ is the leftmost point of its run in $\rho\setminus z_1=\tau+z_2$ if and only if its left neighbour there, the point of $\tau$ at position $i_2-1$, is not adjacent to it in value, i.e. $i_2=1$ or $\tau(i_2-1)\notin\{v_2-1,v_2\}$; likewise for $z_1$ in $\tau+z_1$. So both points must avoid the $2a$ \emph{bad} slots $(i,v)$ with $i\ge2$ and $v\in\{\tau(i-1),\tau(i-1)+1\}$, and there are $a^2+1$ good slots: $a+1$ in position gap $1$ and $a-1$ in each other position gap; $a$ in value gaps $1$ and $a+1$ and $a-1$ in each other value gap. 

A pair of good slots differing in both coordinates gives one edge, a pair differing in exactly one coordinate gives two (the order in the shared gap), and a slot taken twice gives two; hence the number of locally least edges summed over $\rho$ is
\[
\tfrac12\bigl[(a^2+1)^2-(a^2+1)-N_1\bigr]+N_1+2(a^2+1),
\]
where $N_1=a(a+1)+a(a-1)(a-2)+2a(a-1)+(a-1)^2(a-2)$ is the number of ordered pairs of distinct good slots sharing a coordinate; this simplifies to $(a^4+2a^3+a^2+6a+2)/2=C_2(a)+a-1$.
The same polynomial is the number of active two-point insertion triples
in Ray and West \cite[Prop.~5.3, p.~68]{RayWest2003}; the slot count here
uses the leftmost-run normalization. 

Since every component has exactly one least edge, which is locally least, \cref{thm:A} is equivalent to the statement that, summed over $\rho$, exactly $a-1$ locally least edges fail to be least in their component. This count is universal, but the $a-1$ configurations are not uniformly described, which is why we argued through the cycle space.
\end{remark}

\begin{remark}[verification]
The program \texttt{delgraph.cpp} in the companion repository checks, for every $\tau$ of length at most $7$ and every $\rho\in\Sym_{|\tau|+2}$, identity \eqref{eq:b1}, the conclusion of \cref{lem:leading} for every edge closing a cycle, and the sums $\sum_\rho\bz=C_2(a)$ and $\sum_\rho\bone=2a(a+2)$. These are not needed for the proof; they were used to find the statements.
\end{remark}

\section{The interval formula for \texorpdfstring{$j$}{j}}\label{sec:interval}

\Cref{cor:j} evaluates $j(\tau)$ on the permutations of length $a+2$. This section evaluates it on $\tau$ itself: $j(\tau)$ is a sum over the intervals of $\tau$ of a $\{0,1,2\}$-valued weight that can be read off the interval in linear time (\cref{thm:interval,thm:catalogue}). The formula is derived from the codimension-two classification of Ray and West; its consequences -- a nonnegative deficit on the substitution tree, the characterization of the maximizers and of the next-to-maximal patterns, and a Poisson limit law -- are proved from it without any further enumeration.

\subsection{Returns and the statement}

\begin{definition}[return test, primitive returns]\label{def:return}
Let $w=w_1\cdots w_l\in\Sym_l$, $l\ge2$, and put $a:=w_1$ and $b:=w^{-1}(1)$. If $a=1$ or $b=1$ put $R(w):=0$. Otherwise call an index $q$ \emph{eligible} if $b-1\le q\le l-1$, $\max(w_q,w_l)=l$ and $v:=\min(w_q,w_l)\ge a-1$; for an eligible $q$ let
\[
\psi(x):=x+1+[x\ge a-1]-[x>v]
\]
and let $D_q(w)$ be the word of length $l$
\[
D_q(w):=\bigl(a,\ \psi(w_1),\dots,\psi(w_{b-2}),\ 1,\ \psi(w_{b-1}),\dots,\psi(w_{l-1})\bigr)\quad\text{with the entry }\psi(w_q)\text{ omitted}.
\]
Put $R(w):=1$ if $D_q(w)=w$ for some eligible $q$ (\emph{$w$ returns}), and $R(w):=0$ otherwise. The \emph{primitive-return indicator} is
\[
P(w):=R(w)\prod_{t}\bigl(1-R(w_1\cdots w_t)\bigr),
\]
the product over $2\le t<l$ with $\{w_1,\dots,w_t\}=[t]$ (an empty product is $1$): $w$ returns and no proper prefix that is itself a permutation returns.
\end{definition}

The test has a geometric meaning. When $a,b\ge2$, let $\rho_w\in\Sym_{l+2}$ be obtained from $w$ by inserting a point of final value $a$ at the front and a point of final value $1$ immediately before the $(b-1)$st old point:
\[
\rho_w=\bigl(a,\ \phi(w_1),\dots,\phi(w_{b-2}),\ 1,\ \phi(w_{b-1}),\dots,\phi(w_l)\bigr),\qquad \phi(x):=x+1+[x\ge a-1].
\]
In one-based (value gap, position gap) coordinates relative to $w$, the insertion sites are $(a-1,1)$ and $(1,b-1)$. The map $\phi$ records the simultaneous shifts of the old values. The old points form an occurrence $O$ of $w$ in $\rho_w$. The condition $D_q(w)=w$ says that the points of $\rho_w$ other than $\phi(w_q)$ and $\phi(w_l)$ form a second occurrence $O'$ of $w$; one of these two deleted points is the maximum of $\rho_w$, because $\max(w_q,w_l)=l$, and the other is its last point or, when $w_l=l$, the point $\phi(w_q)$. The deletion pairs of $O$ and $O'$ are disjoint. So $R(w)=1$ says that $\rho_w$ contains $w$ in two ways with disjoint deletion pairs, the second lying to the left of and below the first, and $P(w)$ requires in addition that this does not already happen inside a proper initial segment of $w$ carrying the values $[t]$.

\begin{example}
We have $R(21)=P(21)=1$ ($\rho_{21}=2143$, $q=1$); $R(12)=0$ ($a=1$); $R(213)=1$ but $P(213)=0$ because the prefix $21$ returns; $P(231)=P(312)=1$; $P(2413)=P(3142)=1$.
\end{example}

For $w=2413$ we have $a=2$, $b=3$, and the eligible indices $q\in\{2,3\}$ must satisfy $\max(w_q,w_4)=4$, so $q=2$ and $v=\min(4,3)=3\ge a-1$; then $\psi(x)=x+2-[x>3]$ and $D_2(w)=(2,\psi(2),1,\psi(1))=(2,4,1,3)=w$, so $w$ returns, and $P(w)=1$ because $w$ has no proper prefix with value set $[t]$. 

As illustrated in \cref{fig:return}, $\rho_w=241635$ contains $2413$ at the positions $\{2,4,5,6\}$ (the old points, values $4,6,3,5$) and at $\{1,2,3,5\}$ (values $2,4,1,3$), the second occurrence being obtained by deleting the maximum $6$ and the last point $5$; the two deletion pairs $\{1,3\}$ and $\{4,6\}$ are disjoint.

\begin{figure}[htbp]
\centering
\begin{tikzpicture}[x=.48cm,y=.48cm, every node/.style={font=\footnotesize}]
\foreach \shift/\caption in {0/{$O=\{2,4,5,6\}$},11/{$O'=\{1,2,3,5\}$}} {
  \begin{scope}[xshift=\shift*.48cm]
  \draw[step=1,black!18,very thin] (1,1) grid (6,6);
  \draw[->] (.5,.5)--(7,.5) node[right]{position};
  \draw[->] (.5,.5)--(.5,7) node[above]{value};
  \foreach \i in {1,...,6} {\node[below] at (\i,.5) {\i}; \node[left] at (.5,\i) {\i};}
  \foreach \x/\y in {2/4,4/6,5/3,6/5}
    \draw[fill=white,line width=.6pt] (\x,\y) circle (3pt);
  \foreach \x/\y in {1/2,3/1}
    \draw[fill=white,line width=.6pt] (\x-.17,\y-.17) rectangle (\x+.17,\y+.17);
  \node[anchor=north] at (3.5,-.6) {\caption};
  \end{scope}
}
\foreach \x/\y in {2/4,4/6,5/3,6/5} \fill (\x,\y) circle (2pt);
\begin{scope}[xshift=5.28cm]
\foreach \x/\y in {1/2,2/4,3/1,5/3} \fill (\x,\y) circle (2pt);
\end{scope}
\end{tikzpicture}
\caption{The return $2413\mapsto\rho_w=241635$. Circles are old
points and squares are the two inserted points. Filled centers select
the indicated occurrence; its complement is the deletion pair.
Thus $O$ deletes positions $\{1,3\}$, while $O'$ deletes $\{4,6\}$.
Both selected patterns standardize to $2413$.}
\label{fig:return}
\end{figure}

\begin{theorem}[interval formula]\label{thm:interval}
For every $\tau\in\Sym_a$,
\begin{equation}\label{eq:interval}
j(\tau)=\sum_{I}\bigl[P(\std I)+P((\std I)^{r})\bigr],
\end{equation}
the sum over all intervals $I$ of $\tau$ of size at least two, $\tau$ itself included, where $\std I$ is the pattern of $\tau$ on $I$ and $(\cdot)^r$ is reversal.
\end{theorem}

We use the insertion conventions of Ray and West
\cite[\S4, p.~64; \S5, p.~67]{RayWest2003}. A two-point insertion into
$\tau$ is specified by a triple $(\pi;R,C)$, where $\pi\in\Sym_2$ orders
the two inserted points, $R=(r_1\le r_2)$ lists their value gaps, and
$C=(c_1\le c_2)$ lists their position gaps, each numbered from $1$ to
$a+1$. Its insertion sites are $(r_{\pi(i)},c_i)$ for $i=1,2$.
Two triples are \emph{synonymous} when their resulting permutations agree.
For distinct row multisets, the \emph{left-hand} triple is the one with
lexicographically smaller row multiset. A site $(r,c)$ is \emph{active}
when the permutation matrix of $\tau$ has zeros at both $(r,c-1)$ and
$(r,c)$; entries outside the matrix are read as zero. A triple is active
when both sites are active. Ray and West's underlying permutation
$\pi\in\Sym_m$ is our $\tau\in\Sym_a$, their $t_2(\pi)$ is our
$g_2(\tau)$. Their $\eta_m$ denotes the decreasing permutation
\cite[p.~58]{RayWest2003}, so $\eta_2=21$ and $\iota_2=12$.
The inserted order is not restricted in general: Proposition~8.5 forces
$21$ for the unreflected active type-C source, while horizontal
reflection gives source order $12$.
Their type-C classification supplies the following input.

\begin{theorem}[Ray--West {\cite[pp.~80--83]{RayWest2003}}]\label{thm:RW}
For every $\tau\in\Sym_a$, $g_2(\tau)=C_2(a)-j(\tau)$ where $j(\tau)$ is the number of active left-hand insertion triples admitting a synonymity of type C, in either of its two orientations. In the unreflected orientation, a type-C synonymity between a source $(21;(r_1,r_2),(c_1,c_2))$ and a target $(\pi';(s_1,s_2),(d_1,d_2))$ has
\begin{equation}\label{eq:typeC}
r_1\le r_2<s_1\le s_2,\qquad c_1\le c_2<d_1\le d_2,
\end{equation}
and the points of $\tau$ affected by it form a contiguous square submatrix, that is, an interval of $\tau$; the reflected orientation is its horizontal mirror image, with source order $12$. Here $j$ counts sources: several targets may be synonymous with the same source.
\end{theorem}

The source of each part of \cref{thm:RW} is explicit. On p.~80,
Ray and West give the strict inequalities \eqref{eq:typeC} and state that
``the track matrix is a contiguous submatrix of $\pi$'';
Proposition~8.5 on that page specifies the source orientation.
Corollary~8.6 on p.~82 counts the active left-hand triples of type~C,
including the reflected orientation, and Theorem~8.7 on p.~83 gives the
formula for $g_2$. Thus the count is of sources, not of all source--target
pairs. This is the imported classification used in this section; the
other external inputs are listed in \cref{sec:boundary}.
The four-parameter track family in \cref{thm:catalogue} is a
reparametrization of Ray--West Proposition~8.5(1); the coordinate
dictionary is given there. Our additions are the selection of the unique
primitive support, its enumeration, the evaluation on intervals of
$\tau$, and the consequences in \cref{sec:deficit,sec:poisson}.

\begin{proof}[Proof of \cref{thm:interval}]
We show that the unreflected type-C sources are in bijection with the intervals $I$ with $P(\std I)=1$, and the reflected ones with the intervals with $P((\std I)^r)=1$; \cref{thm:RW} then gives \eqref{eq:interval}.

\emph{Normalization.} Let a type-C synonymity in the unreflected orientation have support interval $I$ of size $l$, and let $w=\std I\in\Sym_l$. In coordinates relative to $w$, the extreme gaps are $r_1=c_1=1$ and $s_2=d_2=l+1$. The target deletes two old points from the source expansion; by \eqref{eq:typeC} both target rows lie strictly above $r_2$ and both target columns strictly to the right of $c_2$. Hence the first entry of the target copy of $w$ is the new point of value $r_2+1$, so $r_2=a-1$ with $a=w_1$; and the new point of value $1$ is not deleted and remains at position $c_2+1$ in the target copy, so $c_2=b-1$ with $b=w^{-1}(1)$. Thus the source is $(21;(1,a-1),(1,b-1))$, and its expansion is $\rho_w$.

\emph{Targets.} Since $d_2=l+1$, one deleted point is the last point $\phi(w_l)$ of $\rho_w$; the other is $\phi(w_q)$ with $b-1\le q\le l-1$, because $d_1>c_2=b-1$. Since $s_2=l+1$, the largest value $l+2=\phi(l)$ is among the two deleted values, i.e. $\max(w_q,w_l)=l$. With $v=\min(w_q,w_l)$, the strict row separation $s_1>r_2$ is equivalent to $v\ge a-1$: if $v\le a-2$ then $\phi(v)\le a-1=r_2$, and if $v\ge a-1$ then $\phi(v)\ge a+1>r_2$. The values $a$ and $1$ of the new points lie below both deleted values and are unchanged by standardization, and a retained old value $x$ becomes $\phi(x)-[x>v]=\psi(x)$. Hence the standardization of $\rho_w$ minus the two deleted points is exactly $D_q(w)$, and the source admits a type-C target on the support $I$ if and only if $R(w)=1$; conversely, an eligible $q$ with $D_q(w)=w$ exhibits the target, with gap inequalities \eqref{eq:typeC}.

\emph{Activity and extension to $\tau$.} The source is active: its site in row $1$ and column $b-1$ is not adjacent to the old minimum, which lies in column $b$, and its site in row $a-1$ and column $1$ has $w_1=a\ne a-1$. A site on the boundary of an interval $I$ of $\tau$ stays active in $\tau$, because no point outside $I$ has a value inside the value range of $I$. Conversely, a return inside $I$ extends to a synonymity of $\tau$: both descriptions insert the same two rows and columns inside the boundaries of $I$, and the points outside $I$ are unaffected.

\emph{Each source once.} A fixed unreflected source has a fixed minimum row $r_1$ and minimum column $c_1$, so all its type-C supports are square intervals with the same south-west corner and are nested. Within a support $w$, the smaller supports of the same source are exactly the prefixes $w_1\cdots w_t$ with value set $[t]$ that return; they have the same first value $a$ and the same position $b$ of the minimum, hence the same source triple. Every source admitting a support therefore has a unique smallest support, and $P(w)$ is the indicator that $I$ is that smallest support. The reflected orientation is the same argument for $w^r$; the two orientations have source orders $21$ and $12$ and never describe the same triple. So \eqref{eq:interval} counts every active type-C source exactly once.
\end{proof}

\begin{remark}\label{rem:certB}
The theorem was checked against exact values of $g_2$ for all $46{,}232$ patterns of length $2$ to $8$, and against $366$ further patterns of lengths $9$ to $20$ counted by an independent enumerator (the companion repository, and \cref{sec:S8} for the cluster values).
\end{remark}

\subsection{Structure of returning words}

Matching equal roles in two occurrences gives the chain graph of
Bevan, Homberger and Tenner \cite[preprint, Def.~2.3 and Prop.~2.7]{BHT2018}.
For a return, deleting the two newly inserted points from its two
nontrivial chains gives the paths below. They are also the two tracks
of Ray and West; the lemma records the precise increments in our
coordinates.

\begin{lemma}[two increasing paths]\label{lem:paths}
Let $w\in\Sym_l$ return with eligible index $q$. Define a directed graph on the positions $[l]$ with edges $i\to i+1$ for $1\le i\le b-2$, $i\to i+2$ for $b-1\le i\le q-1$, and $i\to i+1$ for $q+1\le i\le l-1$. Then the graph is the disjoint union of two directed paths, starting at the positions $1$ and $b$ and ending at $q$ and $l$, and $w$ increases along every edge by $1+[a-1\le w_i<v]$. In particular a returning word avoids $321$.
\end{lemma}

\begin{proof}
The edges align each retained old entry with the position it occupies in $D_q(w)$: the entry $w_i$ sits at position $i+1$ of $D_q(w)$ for $i\le b-2$, at position $i+2$ for $b-1\le i<q$, and at position $i+1$ for $i>q$. Every position other than $1$ and $b$ has exactly one incoming edge, and every position other than $q$ and $l$ has exactly one outgoing edge, so the graph consists of two paths with the stated ends. The equation $D_q(w)=w$ says $w_{i^+}=\psi(w_i)$ for every edge $i\to i^+$, and $\psi(x)-x=1+[x\ge a-1]-[x>v]$ equals $1+[a-1\le x<v]$ on the retained values, which are different from $v$ and $l$. Three points forming $321$ cannot lie on two increasing sequences.
\end{proof}

\begin{lemma}[localization at direct-sum cuts]\label{lem:localize}
If $w$ returns and $w=u\oplus u'$ with $|u|=t<l$, then $u$ returns. Consequently
\begin{equation}\label{eq:Pchar}
P(w)=1\iff R(w)=1\text{ and }w\text{ is sum-indecomposable}.
\end{equation}
\end{lemma}

\begin{proof}
Since $w=u\oplus u'$, the first $t$ entries of $w$ are the values $[t]$, so the source parameters $a=w_1$ and $b=w^{-1}(1)$ satisfy $a,b\le t$, and $D_q(w)=w$ says that the first $t$ entries of $D_q(w)$ are also the values $[t]$. By the definition of $D_q(w)$, its first $t$ entries are $a$, $1$ and the images under $\psi$ of the entries of $u$ other than two \emph{locally deleted} ones: if $q<t$ these are $w_q$ and $w_t$ (the entry $\psi(w_q)$ is omitted, and the block of $t-2$ images of $u$-entries then ends with $\psi(w_{t-1})$, so $w_t$ is not among the first $t$ entries either); if $q\ge t$ they are $w_{t-1}$ and $w_t$ (nothing is omitted before position $t$, and the two new entries push the last two entries of $u$ out of the first $t$ places). Put $q':=q$ in the first case and $q':=t-1$ in the second; in both cases $q'\ge b-1$. The two locally deleted values include $t$, because $\psi(t)>t$ and no retained value may leave $[t]$; the other one, $v'$, is determined case by case: if $q<t$ and $w_q<t$ then $w_t=t$ and $v'=w_q=v$; if $q<t$ and $w_q=t$ then $w_t=t-1=v'$; if $q=t$ and $w_t<t$ then $w_{t-1}=t$ and $v'=w_t=v$; if $q=t$ and $w_t=t$ then $v'=t-1$; if $q>t$ both globally deleted values exceed $t$, and then $t$ and $t-1$ are exactly the two values of $[t]$ whose images under $x\mapsto x+1+[x\ge a-1]$ exceed $t$. In every case $v'\ge a-1$, and on the retained values of $u$ the local map $x\mapsto x+1+[x\ge a-1]-[x>v']$ agrees with $\psi$. So the first $t$ entries of $D_q(w)=w$ form the return equation of $u$ with deleted positions $q',t$, and $u$ returns.

A returning word has no initial component of size one, since $a,b\ge2$. Hence a proper direct-sum cut of a returning $w$ gives a proper returning prefix with value set $[t]$, and $P(w)=0$; conversely a sum-indecomposable word has no prefix with value set $[t]$, $t<l$, so $P(w)=R(w)$.
\end{proof}

\begin{theorem}[primitive selection in the Ray--West catalogue]\label{thm:catalogue}
Let $a,b\ge2$ and $t,s\ge0$ be integers with
\begin{equation}\label{eq:cat}
l:=a+b+t+s-2,\qquad q:=b-1+2t<l,\qquad v:=a-1+2s<l .
\end{equation}
Put $w_1:=a$, $w_b:=1$, and propagate along the edges of \cref{lem:paths} (with these $b$, $q$, $l$) by $x\mapsto x+1+[a-1\le x<v]$. This defines a permutation $Z(a,b;t,s)\in\Sym_l$, and the words with $P(w)=1$ are exactly the $Z(a,b;t,s)$ except those with $a=b=2$ and $t=s\ge1$, which are $21\oplus21\oplus\dots\oplus21$. The parameters are recovered from $w$ as $a=w_1$, $b=w^{-1}(1)$, $q=w^{-1}(l)$, $v=w_l$, $t=(q-b+1)/2$, $s=(v-a+1)/2$; in particular $P(w)$ can be decided in $O(l)$ steps.
\end{theorem}

\begin{proof}
Let $P(w)=1$. Then $w_l\ne l$: otherwise $w=u\oplus1$ is decomposable, contradicting \eqref{eq:Pchar}. So $q=w^{-1}(l)$ and $v=w_l$ are forced. By \cref{lem:paths} the positions form two increasing paths starting at $1,b$ and ending at $q,l$; correspondingly the values form two chains starting at $a,1$ (the roots) and ending at $l,v$ (the deleted values), each value chain being the image of a position path. 

Put $k:=q-b+1$ and $m:=v-a+1$. The path starting at position $1$ ends at $q$ if and only if $k$ is even (the path from $1$ passes $b-1$ and then takes steps of $2$ until $q$ or $q-1$), and the value chain from $a$ ends at $l$ if and only if $m$ is even (it takes steps of $2$ between $a-1$ and $v$). Since $w_q=l$ and $w_l=v$, the path ending at $q$ carries the chain ending at $l$, so $k$ and $m$ have the same parity. If $k=2t+1$ and $m=2s+1$ were odd, the position path from $1$ would have $l-t-1$ vertices and the value chain from $a$ would have $s+1$, forcing $l=t+s+2$; but $q<l$ and $v<l$ give $t-s\le1-b$ and $s-t\le1-a$, whose sum contradicts $a,b\ge2$. So $k=2t$, $m=2s$, the path from $1$ has $b-1+t$ vertices and the chain from $a$ has $l-a+1-s$, and equating gives \eqref{eq:cat}. Conversely, under \eqref{eq:cat} the paired paths and chains have equal lengths, so the construction defines a bijection $[l]\to[l]$ satisfying the return equation with eligible index $q$; so $Z(a,b;t,s)$ returns. 

\emph{Primitivity.} It remains to identify the decomposable constructions, i.e.\ the returning words $Z=Z(a,b;t,s)$ with $P(Z)=0$. Suppose $Z=U\oplus U'$ with $U$ sum-indecomposable of size $k<l$. Then $a=w_1\le k$ and $b=w^{-1}(1)\le k$, while $q=w^{-1}(l)>k$ and $v=w_l>k$ because $l$ and $w_l$ belong to $U'$. We are therefore in the case $q>t=k$ of \cref{lem:localize}: $U$ returns with its last two positions $k-1,k$ locally deleted, and the locally deleted values are $k$ and $k-1$. So $\{U_{k-1},U_k\}=\{k-1,k\}$, and $U_k=k$ would make $U=U''\oplus1$ decomposable; hence $U_k=k-1$ and $U_{k-1}=k$. Now $U$ is sum-indecomposable and returns, so $P(U)=1$ by \eqref{eq:Pchar} and, by the first part of the proof applied to $U$, $U=Z(a,b;t_U,s_U)$ with $q_U=k-1=b-1+2t_U$ and $v_U=U_k=k-1=a-1+2s_U$, and $k=a+b+t_U+s_U-2$. Substituting $t_U=(k-b)/2$ and $s_U=(k-a)/2$ gives $k=a+b+k-\frac{a+b}2-2$, i.e.\ $a+b=4$, so $a=b=2$. With $a=b=2$ the constraints $q=1+2t<l=2+t+s$ and $v=1+2s<l$ read $t\le s$ and $s\le t$, so $t=s$, and propagating from $w_1=2$, $w_2=1$ (both paths advance by two positions and by two values until the deleted values are reached) gives $Z(2,2;t,t)=21\oplus\dots\oplus21$ with $t+1$ blocks. Conversely, every $21\oplus\dots\oplus21$ with at least two blocks is decomposable and returns ($q=l-1$, $v=l-1$), so it is exactly the excluded family. Uniqueness of the parameters is clear from their definition.
\end{proof}

\paragraph{Comparison with the Ray--West catalogue.}
In Proposition~8.5(1) of \cite[pp.~80--81]{RayWest2003}, let
$h_1,h_2,h_3$ and $w_1,w_2,w_3$ be the successive row and column
block sizes. For $Z(a,b;t,s)$ the dictionary is
\[
\begin{array}{lll}
h_1=a-2,&h_2=2(s+1),&h_3=l-a-2s,\\
w_1=b-2,&w_2=2(t+1),&w_3=l-b-2t.
\end{array}
\]
Indeed the source gaps are $(1,a-1)$ and $(1,b-1)$, and the target
gaps in this form are $(v+2,l+1)$ and $(q+2,l+1)$.
Taking consecutive differences gives the displayed block sizes.
With Ray--West's $h=h_2/2=s+1$ and $w=w_2/2=t+1$, their condition
$h_1-w=w_3-h$ is exactly $l=a+b+t+s-2$.
Their second form and its transpose have target order $12$, so the last
point of the support is its maximum; these returns have the form
$u\oplus1$ and are nonprimitive by \cref{lem:localize}.
Within the first form, our primitivity condition removes the repeated
$21$ family. Thus the track classification is inherited; the
minimal-support selection and the count below are the additional steps.
For example, their track matrix $61823457$ is $Z(6,2;1,1)$.

\begin{corollary}[the count of primitive returns]\label{cor:rl}
The number $r_l$ of $w\in\Sym_l$ with $P(w)=1$ is
\[
r_l=\frac{l^3-l}{12}\ (l\text{ odd}),\qquad r_l=\frac{l^3+2l}{12}-[l\ge4]\ (l\text{ even}),\qquad
\sum_{l\ge2}r_lx^l=\frac{x^2(1+x^2)}{(1-x)^2(1-x^2)^2}-\frac{x^4}{1-x^2},
\]
so $r_l=1,2,5,10,18,28,43,60,84,110,145$ for $l=2,\dots,12$. Moreover
\begin{equation}\label{eq:mass}
\sum_{\tau\in\Sym_m}j(\tau)=2\sum_{l=2}^{m}r_l\,(m-l+1)^2(m-l)!.
\end{equation}
\end{corollary}

\begin{proof}
Write $d=t-s$ and $u=\min(t,s)$; the constraints in \eqref{eq:cat} become $a-2\ge\max(d,0)$ and $b-2\ge\max(-d,0)$, and after subtracting these minima the remaining parameters $a_0,b_0,u\ge0$ and $d\in\ZZ$ satisfy $l-2=a_0+b_0+2u+2|d|$; the generating function of these quadruples is the first term. The excluded words contribute one at every even $l\ge4$, the second term. Extracting coefficients gives the closed forms. For \eqref{eq:mass}: an interval of size $l$ with a given pattern occurs at $m-l+1$ positions with $m-l+1$ value ranges, the remaining $m-l$ values being arbitrary, and both orientations are counted by \eqref{eq:interval}.
\end{proof}

\begin{remark}[pure returns]\label{rem:pure}
Call $\rho\in\Sym_{l+2}$ a \emph{pure return} of $w$ if $\rho$ contains exactly two occurrences of $w$, with disjoint deletion pairs, one lying to the right of and above the other point by point. Exhaustive comparison shows that $P(w)=1$ if and only if $w$ has a pure return in which the upper-right occurrence has the two new points as its private points, for all $w$ of length at most $7$, and \cref{thm:interval} with this substitute definition is verified for all patterns of length at most $8$ (the recursive check of \cref{sec:S8}). This is the language of \cref{cor:j}: the two occurrences of a pure return lie in different one-move classes. We have not derived this equivalence from \cref{thm:RW}, and do not use it.
\end{remark}

A direct derivation of \cref{thm:interval} from \cref{cor:j} would avoid
importing the insertion classification. The deletion graphs offer a
possible route: exhaustive checks through pattern length seven show that
distinct non-isolated components are separated in both position and
value. Four components already occur for $\tau=214365$ in
$\rho=21436587$, with vertex sets $\{1,2\},\{3,4\},\{5,6\},\{7,8\}$.
This remains finite evidence, not a proved separation theorem; the chain
graphs of \cite[preprint, \S2.1]{BHT2018} provide a natural language for
studying the comparison.

\subsection{The deficit on the substitution tree}\label{sec:deficit}

Put $D(\tau):=|\tau|-1-j(\tau)$, the \emph{deficit}. A permutation is \emph{layered} if it is a direct sum of decreasing permutations and \emph{reverse-layered} if it is the reversal of a layered permutation, i.e. a skew sum of increasing ones.

\begin{theorem}[boundary law]\label{thm:boundary}
For nonempty $\alpha,\beta$,
\begin{equation}\label{eq:boundary}
j(\alpha\oplus\beta)=j(\alpha)+j(\beta)+t(\alpha)\,s(\beta),
\end{equation}
where $t(\alpha)=1$ if the last sum-indecomposable component of $\alpha$ is decreasing and $s(\beta)=1$ if the first sum-indecomposable component of $\beta$ is decreasing, and $t,s=0$ otherwise. Consequently, if $\tau=\tau_1\oplus\dots\oplus\tau_k$ with sum-indecomposable summands,
\begin{equation}\label{eq:sumlaw}
j(\tau)=\sum_i j(\tau_i)+\#\{i<k:\tau_i,\tau_{i+1}\text{ both decreasing}\},
\end{equation}
and dually for skew sums with ``increasing''.
\end{theorem}

\begin{proof}
In \eqref{eq:interval} the intervals inside $\alpha$ or inside $\beta$ contribute $j(\alpha)+j(\beta)$. Let $I$ straddle the cut. Its pattern $w$ is sum-decomposable, so $P(w)=0$ by \eqref{eq:Pchar}. Its reversal $w^r$ is skew-decomposable; if $w^r$ returns it avoids $321$ (\cref{lem:paths}), and a skew-decomposable $321$-avoider has exactly two skew components, both increasing (an inversion inside a component together with a point of the other, or three components, gives $321$). So $w=\dd_r\oplus\dd_s$ with $r,s\ge1$. Conversely $w^r=(r+1)(r+2)\cdots(r+s)\,12\cdots r$ is sum-indecomposable and returns (delete its last point and the last point of its first block: $a=r+1$, $b=s+1$, $q=s$, $v=r$), so $P(w^r)=1$. A straddling interval of the form $\dd_r\oplus\dd_s$ consists of a decreasing suffix of $\alpha$ carrying the $r$ largest values of $\alpha$ and a decreasing prefix of $\beta$ carrying the $s$ smallest values of $\beta$; such a suffix is a union of final components of $\alpha$, and being decreasing it is a single decreasing final component, and likewise for $\beta$. Hence there is exactly one contributing straddling interval if $t(\alpha)=s(\beta)=1$ and none otherwise. Iterating gives \eqref{eq:sumlaw}; reversal gives the skew version.
\end{proof}

For a decreasing $\dd_r$ only the $r-1$ adjacent intervals contribute (a decreasing word of length $\ge3$ contains $321$, and its reversal begins with $1$), so $j(\dd_r)=r-1$ and, by \eqref{eq:sumlaw}, $j(\tau)=|\tau|-1$ for every layered $\tau$. In particular $g_2(\tau)=(a^4+2a^3+a^2+2a+6)/2$ for layered $\tau$ of length $a$.

\begin{lemma}[simple inflations]\label{lem:simple}
If $\tau=\sigma[\tau_1,\dots,\tau_r]$ with $\sigma$ simple of length $r\ge4$, then $j(\tau)=\sum_i j(\tau_i)+P(\tau)+P(\tau^r)$.
\end{lemma}

\begin{proof}
A proper interval of $\tau$ meeting two blocks projects to a nontrivial interval of $\sigma$ unless it meets every block; the first and last entries of a simple permutation of length at least four are not extreme, so the blocks carrying the global minimum and maximum are interior, an interval meeting every block contains them entirely and is all of $\tau$. Thus the proper intervals lie in single blocks.
\end{proof}

\begin{corollary}[double weight]\label{cor:double}
$P(w)+P(w^r)\le2$, with equality only for $w\in\{2413,3142\}$.
\end{corollary}

\begin{proof}
If $P(w)=P(w^r)=1$ then $w$ avoids $321$ and, by reversal, $123$. Assigning to each entry the lengths of the longest increasing and decreasing subsequences ending there gives distinct pairs in $\{1,2\}^2$, so $|w|\le4$. By \eqref{eq:Pchar} applied to $w$ and to $w^r$, $w$ is both sum- and skew-indecomposable; the only such permutations of length at most $4$ are $1$, $2413$ and $3142$, and the last two satisfy $P(w)=P(w^r)=1$ by \cref{thm:catalogue} ($2413=Z(2,3;0,1)$, $3142=Z(3,2;1,0)$).
\end{proof}

\begin{theorem}[extremal patterns]\label{thm:extremal}
For every $\tau$ of length $a\ge1$, $0\le j(\tau)\le a-1$, and $j(\tau)=a-1$ if and only if $\tau$ is layered or reverse-layered. For $a\ge2$ exactly $2^a-2$ patterns attain the maximum.
\end{theorem}

\begin{proof}
By the substitution decomposition \cite{BHV2008}, a permutation of length $\ge2$ is a direct sum of sum-indecomposables, a skew sum of skew-indecomposables, or an inflation of a simple permutation of length $\ge4$. By \eqref{eq:sumlaw}, \cref{lem:simple} and \cref{cor:double},
\[
D(\tau)=\sum_i D(\tau_i)+\#\{i<k:\tau_i,\tau_{i+1}\text{ not both decreasing}\}
\]
at a direct-sum node (dually at a skew node), and $D(\tau)=\sum_iD(\tau_i)+r-1-P(\tau)-P(\tau^r)\ge\sum_iD(\tau_i)+r-3$ at a simple node. Every local term is nonnegative and $D(1)=0$, so $D\ge0$, i.e. $j\le a-1$; the lower bound $j\ge0$ is \eqref{eq:interval} (or \cref{cor:j}). A simple root contributes at least $1$. A direct-sum root contributes $0$ only if every adjacent pair of components is a pair of decreasing permutations, i.e. all components are decreasing and $\tau$ is layered; conversely layered permutations have $D=0$ by \eqref{eq:sumlaw}. The skew case gives the reverse-layered permutations. There are $2^{a-1}$ layered permutations of length $a$, as many reverse-layered ones, and the two families meet exactly in $\iota_a$ and $\dd_a$.
\end{proof}

\begin{theorem}[next-to-maximal patterns]\label{thm:nextmax}
$j(\tau)=|\tau|-2$ if and only if exactly one of the following holds: $\tau\in\{2413,3142\}$; the direct-sum decomposition of $\tau$ has at least two components, all layered or reverse-layered, and exactly one adjacent pair of components is not a pair of decreasing permutations; or the same with skew sums and increasing permutations. The number $h_a$ of such patterns of length $a$ has generating function
\begin{equation}\label{eq:H}
\sum_{a\ge0}h_ax^a=\frac{8x^4}{(1-x)(1-2x)^2}+\frac{8x^6}{(1-x)^2(1-2x)^2}+2x^4,
\end{equation}
so $h_a=10,40,144,440,1216,3144,7760$ for $a=4,\dots,10$.
\end{theorem}

\begin{proof}
Deficit one at a direct-sum root requires zero deficit of all components (hence, by \cref{thm:extremal}, layered or reverse-layered components) and exactly one adjacency that is not a pair of decreasing permutations; a simple root has deficit $r-1-P-P^r\ge r-3$, which equals $1$ only for $r=4$ with $P(\tau)=P(\tau^r)=1$, i.e. $\tau\in\{2413,3142\}$ by \cref{cor:double}. For the count, a deficit-zero sum-indecomposable component is decreasing (type A, generating function $x/(1-x)$) or reverse-layered but not monotone (type B: for size $n\ge3$ there are $2^{n-1}-2$ of them, generating function $2x^3/((1-x)(1-2x))$). Exactly one bad adjacency means one type-B component at an end followed or preceded by a nonempty sequence of type-A components, or exactly two components both of type B; this gives $2BA/(1-A)+B^2$ for direct-sum roots, the same for skew roots, and $2x^4$ for the two simple exceptions.
\end{proof}

By \cref{thm:extremal}, no pattern of length $a$ that is neither layered nor reverse-layered is Wilf-equivalent to a layered one: its $g_2$ is larger by at least one. On $\Sym_8$ the numbers of patterns with $j=0,\dots,7$ are $4632$, $8366$, $9550$, $7958$, $5416$, $2928$, $1216$, $254$ (from the data of \cref{sec:S8}), in agreement with $h_8=1216$, $2^8-2=254$ and the total $\sum j=96718$ from \eqref{eq:mass}.

\subsection{A Poisson law}\label{sec:poisson}

\begin{corollary}[Poisson law]\label{cor:poisson}
Let $\Pi_n$ be uniformly distributed in $\Sym_n$ and let $B_n$ be its number of bonds. Then $j(\Pi_n)=B_n+R_n$ with $R_n\ge0$,
\begin{equation}\label{eq:ER}
\EE R_n=\frac{2}{n!}\sum_{l=3}^{n}r_l\,(n-l+1)^2(n-l)!=\frac4n+\frac6{n^2}+O(n^{-3}),
\end{equation}
$\EE\,j(\Pi_n)=2+\frac2n+\frac6{n^2}+O(n^{-3})$, and $j(\Pi_n)$ converges in distribution to a Poisson random variable with mean $2$. (The Poisson law for the number of bonds is classical \cite{Wolfowitz1944,Kaplansky1945}; the content of the corollary is that $j$ differs from it by $R_n$.)
\end{corollary}

\begin{proof}
An interval of size two contributes exactly one to \eqref{eq:interval} ($P(21)=1=P(12^r)$, $P(12)=0=P(21^r)$), so $j=B_n+R_n$ with $R_n$ the contribution of the intervals of size $\ge3$. For fixed positions $i,\dots,i+l-1$ and $\pi\in\Sym_l$, the probability that these positions form an interval of $\Pi_n$ with pattern $\pi$ is $(n-l+1)(n-l)!/n!$; summing over the $n-l+1$ starting positions and using $\sum_{\pi}[P(\pi)+P(\pi^r)]=2r_l$ gives the first equality in \eqref{eq:ER} (this is \eqref{eq:mass} divided by $n!$, without the $l=2$ term). The terms $l=3,4$ are $4(n-2)/(n(n-1))$ and $10(n-3)/(n(n-1)(n-2))$, whose sum is $4/n+6/n^2+O(n^{-3})$. For the tail use $r_l\le l^3/8$ (\cref{cor:rl}): for $5\le l\le n/2$ each term is at most $\tfrac14l^3n^2(2/n)^l$ because $(n-l)!/n!\le(2/n)^l$ there, and these sum to $O(n^{-3})$; for $l>n/2$ the terms are bounded by a polynomial in $n$ times $\lfloor n/2\rfloor!/n!$. With $\EE B_n=2(n-1)/n$ this gives the expansion of $\EE\,j(\Pi_n)$.

For the limit law, $\PP(j(\Pi_n)\ne B_n)\le\EE R_n\to0$, so it suffices to treat $B_n$, and the Poisson law for the total number of bonds is classical
\cite{Wolfowitz1944,Kaplansky1945}. The following moment computation
proves the limit needed here. Identify a bond with the value pair $\{v,v+1\}$ it carries. For $1\le k\le n-1$ choose $k$ of the $n-1$ pairs; if they form $r$ maximal runs of consecutive pairs, the permutations in which all $k$ pairs are adjacent are obtained by contracting each run to a contiguous monotone block in one of two directions and arranging the $n-k$ resulting objects, and there are $\binom{k-1}{r-1}\binom{n-k}r$ choices with $r$ runs. Hence
\[
\EE\binom{B_n}{k}=\frac{(n-k)!}{n!}\sum_{r=1}^{\min(k,n-k)}2^r\binom{k-1}{r-1}\binom{n-k}{r}\xrightarrow[n\to\infty]{}\frac{2^k}{k!},
\]
the binomial moments of the Poisson distribution with mean $2$, which is determined by its moments.
\end{proof}

\section{Codimension three}\label{sec:codim3}

By \eqref{eq:g3cluster} and \cref{thm:universal}, $g_3(\tau)$ has the following equivalent descriptions:
\begin{align}
g_3(\tau)&=U_3(a)-(a+3)^2j(\tau)+c_{a+3}(\tau) && \text{(cluster form)},\label{eq:g3a}\\
&=C_3(a)-(a+3)^2j(\tau)-\delta_3(\tau) && \text{where }c_{a+3}(\tau)=-N_3(a)-\delta_3(\tau),\label{eq:g3b}\\
&=C_3(a)-\Delta_3(\tau) && \text{for }a\le6\text{ (\cref{thm:k34}), conjecturally for all }a,\label{eq:g3c}
\end{align}
where $\delta_3(\tau)=-D_{a+3}(\tau)$ is the deviation of the fourth cluster number from its universal part (\cref{prop:weights}) and $\Delta_3(\tau):=(a+3)^2j(\tau)+\delta_3(\tau)=C_3(a)-g_3(\tau)$. The first two forms are exact identities valid for every $\tau$; $j(\tau)$ is given by \cref{thm:interval}, and $c_{a+3}(\tau)$ is a signed count over hosts of size $a+3$. The third form adds an interpretation: for $a\le6$, $\Delta_3(\tau)$ is the number of extra one-move classes of three-point insertions (\cref{thm:k34}); for larger $a$ this interpretation is conjectural, and nothing below depends on it. This section gives an explicit evaluation of $g_3(\tau)$ for every $\tau$ by a canonical-representative count (\cref{thm:boxes}) and a closed form for layered patterns (\cref{thm:layered}); \cref{thm:separated} adds $g_3(\tau)=C_3(a)$ for every $\tau$ with $d_\infty(\tau)\ge4$.

\subsection{The deletion-overlap formula}\label{sec:boxes}

Fix $\tau\in\Sym_a$. For a set $Q$ of positions write $\delta_Q(\tau):=\std(\tau_i:i\notin Q)$. A permutation of length $a+3$ together with a distinguished occurrence of $\tau$ (a \emph{marked host}) is encoded by
\begin{itemize}
\item $c=(c_1\le c_2\le c_3)$ with $0\le c_j\le a$, where $c_j$ is the number of old points to the left of the $j$th new point (new points numbered from left to right);
\item $\sigma\in\Sym_3$, where $\sigma_j$ is the rank of the value of the $j$th new point among the three new values;
\item $y=(y_1\le y_2\le y_3)$ with $0\le y_i\le a$, where $y_i$ is the number of old values below the $i$th smallest new value;
\end{itemize}
new points sharing a column gap are ordered by index and new points sharing a row gap by $\sigma$-rank. There are $6\binom{a+3}{3}^2$ marked hosts. Let $\mathcal C_a:=\{(c_1\le c_2\le c_3)\}\subseteq\{0,\dots,a\}^3$.

\begin{definition}[deletion boxes]\label{def:boxes}
For $c\in\mathcal C_a$ and $\sigma\in\Sym_3$, the collection $\mathcal B_\tau(c,\sigma)$ of integer boxes in $[0,a]^3$ consists of the boxes obtained as follows. Choose $1\le d\le\min(3,a)$, new indices $X=\{j_1<\dots<j_d\}\subseteq[3]$ and old positions $P=\{p_1<\dots<p_d\}\subseteq[a]$ with
\begin{equation}\label{eq:earlier}
c_{j_1}<p_1 .
\end{equation}
Put $q_r:=c_{j_r}-|P\cap[c_{j_r}]|+r$ for $1\le r\le d$ and $Q:=\{q_1,\dots,q_d\}$; the $q_r$ are increasing and lie in $[a]$. Keep the choice only if
\begin{equation}\label{eq:keep}
\delta_P(\tau)=\delta_Q(\tau)\qquad\text{and}\qquad \std(\sigma_{j_1},\dots,\sigma_{j_d})=\std(\tau_{q_1},\dots,\tau_{q_d}).
\end{equation}
Let $0=v_0<v_1<\dots<v_{a-d}<v_{a-d+1}=a+1$ with $\{v_1,\dots,v_{a-d}\}=\{\tau_i:i\notin P\}$, and $t_r:=\tau_{q_r}-1-\#\{s:\tau_{q_s}<\tau_{q_r}\}$. Start with $L_i=0$, $U_i=a$ ($i=1,2,3$), replace $L_{\sigma_{j_r}}:=v_{t_r}$ and $U_{\sigma_{j_r}}:=v_{t_r+1}-1$ for each $r$, and include the box $[L_1,U_1]\times[L_2,U_2]\times[L_3,U_3]$.
\end{definition}

All tests in the definition compare $\tau$ with its own deletion minors; no permutation longer than $\tau$ occurs.

\begin{theorem}[deletion-overlap formula]\label{thm:boxes}
For every $\tau\in\Sym_a$, $a\ge1$,
\begin{equation}\label{eq:boxes}
g_3(\tau)=6\binom{a+3}{3}^2-\sum_{c\in\mathcal C_a}\sum_{\sigma\in\Sym_3}\#\Bigl\{y=(y_1\le y_2\le y_3):\ y\in\bigcup\mathcal B_\tau(c,\sigma)\Bigr\}.
\end{equation}
\end{theorem}

\begin{proof}
Among the markings of a host containing $\tau$, keep only the one whose occurrence is lexicographically first as an increasing list of positions; then each containing host is counted exactly once, and \eqref{eq:boxes} follows once we show that, for fixed $(c,\sigma)$, the marked host $(c,\sigma,y)$ is \emph{not} the first marking if and only if $y$ lies in a box of $\mathcal B_\tau(c,\sigma)$.

A different occurrence of $\tau$ in the marked host replaces $d$ old points, at positions $P$, by $d$ new points, with indices $X$, for some $1\le d\le\min(3,a)$. The two occurrences first differ either at the first omitted old point or at the first chosen new point, and the new occurrence is earlier exactly when its first chosen new point precedes the first omitted old point, which is \eqref{eq:earlier}. In the new occurrence the new point $j_r$ has $c_{j_r}-|P\cap[c_{j_r}]|$ old points and $r-1$ new points to its left, so it occupies the position $q_r$ of $\tau$. 

The surviving old points must realize $\tau$ with the positions $Q$ removed, and the selected new points must have the mutual order prescribed by $\tau$ on $Q$: these are the two conditions \eqref{eq:keep}. It remains to compare the selected new points with the surviving old points: the new point destined for position $q_r$ must have exactly $t_r$ surviving old values below it; with the surviving old values $v_1<\dots<v_{a-d}$, a new point in row gap $y$ has this property exactly when $v_{t_r}\le y\le v_{t_r+1}-1$, and this is the constraint on the axis $\sigma_{j_r}$ of the box. Unselected new points impose nothing. Hence an earlier occurrence exists if and only if $y$ lies in one of the boxes, every box coming from a specific earlier occurrence and conversely.
\end{proof}

The right-hand side of \eqref{eq:boxes} can be evaluated in $O(a^6)$
arithmetic operations. Precompute the deletion minors with at most three
deletions and assign identical minors identical integer identifiers.
There are $O(a^3)$ such minors, each constructible in $O(a)$ operations;
lexicographic sorting and identification takes $O(a^4\log a)$ operations.
For each of the $O(a^3)$ column triples and six orders of the inserted
points, there are $O(a^3)$ choices in \cref{def:boxes}, and each test
then uses a bounded number of operations and identifier comparisons.
Add the resulting boxes to a three-dimensional difference array and
perform three prefix-sum passes on its $O(a^3)$ entries. Counting the
positive entries with $y_1\le y_2\le y_3$ gives the union size, with
the claimed total bound. The bound counts arithmetic operations; bit complexity is not asserted.
The companion archive contains two implementations. For lengths $4$
and $5$ they agree with the cluster values of $g_3$ for all $144$ patterns,
and for lengths $6,7,8$ with the values for all $91,595,4755$ class
representatives of \cref{sec:S8}.

\subsection{Layered patterns}

\begin{theorem}[layered patterns]\label{thm:layered}
Let $\tau=\dd_{b_1}\oplus\dots\oplus\dd_{b_t}$ be layered of length $a\ge2$. Then
\begin{equation}\label{eq:layered}
g_3(\tau)=M_3(a)-\#\{i:1<i<t,\ b_i\ge2\},\qquad M_3(a):=g_3(\iota_a),
\end{equation}
where $M_3(a)=\tfrac16\bigl(a^6+6a^5+10a^4+4a^3+19a^2+62a+66\bigr)$, and the same holds for reverse-layered $\tau$.
\end{theorem}

The polynomial $M_3(a)$ already appears in
Ekhad, Shar and Zeilberger \cite[p.~2]{ESZ2015}.
Polynomiality of $g_k(\iota_a)$ for fixed $k$ was already obtained by
Ray and West \cite[Prop.~9.6, p.~87]{RayWest2003}. The value $M_3(a)$:
by Schensted's theorem $g_3(\iota_a)=\sum_{\lambda\vdash a+3,\ \lambda_1\ge a}(f^\lambda)^2$, a sum over the seven shapes $(a+3)$, $(a+2,1)$, $(a+1,2)$, $(a+1,1,1)$, $(a,3)$, $(a,2,1)$, $(a,1,1,1)$ (the shape $(a,3)$ is absent for $a=2$, where its hook-length expression vanishes), and the hook-length formula gives the polynomial; equivalently, \cref{cor:LIS} with $n=a+3\le2a+1$ gives
\begin{equation}\label{eq:M3}
M_3(a)=6\binom{a+3}3^2-4a\binom{a+3}2^2+(a+1)(2a+1)(a+3)^2-\binom{2a+4}3 .
\end{equation}
The correction term is proved in \cref{app:layered} by a recurrence for $g_3(\alpha\oplus1)$. The interior layers are what matters: the first and last layers of $\tau$ contribute nothing, and an interior layer contributes one whenever it has size at least two, whatever its size. For instance $g_3(1234)=2279$ and $g_3(1324)=2278$, so $|\Av_7(1234)|=2761$ and $|\Av_7(1324)|=2762$.
This example is already recorded by Ray and West
\cite[p.~88]{RayWest2003}, who credit West's thesis and B\'ona;
the new statement here is the correction for arbitrary layered patterns.

\begin{theorem}[the Wilf class of the increasing pattern]\label{thm:monoclass}
For $m\ge2$ let $\mathcal E_m=\{\dd_m\}\cup\{\dd_r\oplus\iota_k\oplus\dd_s: r,s\ge1,\ k\ge0,\ r+k+s=m\}$, the layered permutations of length $m$ with no interior nonsingleton layer. Then
\[
\{\pi\in\Sym_m:\pi\sim\iota_m\}=\mathcal E_m\cup\mathcal E_m^{r},
\]
a set of $m(m-1)$ permutations, and $\pi\sim\iota_m$ if and only if $g_2(\pi)=g_2(\iota_m)$ and $g_3(\pi)=g_3(\iota_m)$. For $m\ge4$ the condition on $g_3$ cannot be dropped.
\end{theorem}

\begin{proof}
If $g_2(\pi)=g_2(\iota_m)$ then $j(\pi)=m-1$, so $\pi$ is layered or reverse-layered by \cref{thm:extremal}, and if also $g_3(\pi)=g_3(\iota_m)$ then \eqref{eq:layered} forces every interior layer to be a singleton; so $\pi\in\mathcal E_m\cup\mathcal E_m^r$. Conversely, for $\pi=\dd_r\oplus\iota_k\oplus\dd_s$ the Backelin--West--Xin equivalence \eqref{eq:BWX} gives $\pi\sim\iota_{r+k}\oplus\dd_s$; reverse-complement turns this into $\dd_s\oplus\iota_{r+k}$, and one more application of \eqref{eq:BWX} gives $\iota_m$. The case $\dd_m$ is \eqref{eq:BWX} itself, and reversal handles $\mathcal E_m^r$. 

\emph{Cardinality.} $|\mathcal E_m|=1+\binom m2$ (one one-layer element and $\binom m2$ compositions $(r,1^k,s)$), and $\mathcal E_m\cap\mathcal E_m^r=\{\iota_m,\dd_m\}$, so the union has $2(1+\binom m2)-2=m(m-1)$ elements. For $m\ge4$ the layered pattern $\iota_1\oplus\dd_2\oplus\iota_{m-3}$ has one interior nonsingleton layer: its $g_2$ equals $g_2(\iota_m)$ and its $g_3$ is one less.
\end{proof}

For $m=8$ this is the $56$-element class of $12345678$ in \cref{thm:E-intro}, and the certificate of \cref{sec:S8} confirms it.

\section{The Wilf classification of \texorpdfstring{$\Sym_8$}{S8}}\label{sec:S8}

\subsection{Known equivalences}\label{sec:known}
The eight symmetries of the square preserve Wilf classes. Backelin, West and Xin \cite{BWX2007} constructed, for every $k\ge1$ and every $\gamma$, a bijection between $\Av_n(\iota_k\oplus\gamma)$ and $\Av_n(\dd_k\oplus\gamma)$; thus
\begin{equation}\label{eq:BWX}
\iota_k\oplus\gamma\sim\dd_k\oplus\gamma\qquad\text{for all }k\ge1\text{ and all }\gamma .
\end{equation}
Stankova and West \cite{StankovaWest2002} proved that the permutations $(m-1,m-2,m,\tau)$ and $(m-2,m,m-1,\tau)$ are Wilf-equivalent for every $\tau\in\Sym_{m-3}$; taking complements,
\begin{equation}\label{eq:SW}
231\oplus\gamma\sim312\oplus\gamma\qquad\text{for all }\gamma .
\end{equation}
The case $k=3$ of \eqref{eq:BWX} is due to Babson and West
\cite{BabsonWest2000}. Krattenthaler \cite{Krattenthaler2006} gives a
growth-diagram proof of the Backelin--West--Xin equivalence
\eqref{eq:BWX}. Bloom and Saracino \cite{BloomSaracino2014} give a
bijection between $231$-avoiding and $312$-avoiding placements,
providing a bijective proof of the shape-Wilf equivalence underlying
\eqref{eq:SW}. In the language of the survey \cite{Vatter2026}, \eqref{eq:BWX} and \eqref{eq:SW} say that $\iota_k$ and $\dd_k$, and $231$ and $312$, are shape-Wilf-equivalent, and shape-Wilf equivalence is preserved by right direct sums; only the ordinary statements \eqref{eq:BWX} and \eqref{eq:SW} are used below. For $m=4$ there is one further equivalence, $1342\sim2413$ \cite{Stankova1994}, which is not of this form.

Let $\equiv_m$ be the equivalence relation on $\Sym_m$ generated by the eight symmetries, the moves $\iota_k\oplus\gamma\leftrightarrow\dd_k\oplus\gamma$ ($k\ge2$) and the moves $231\oplus\gamma\leftrightarrow312\oplus\gamma$. Every $\equiv_m$-class is contained in a Wilf class.

\begin{proposition}\label{prop:classes}
The numbers of $\equiv_m$-classes for $m=4,5,6,7,8$ are $4$, $16$, $91$, $595$ and $4{,}755$. Adding the move $1342\leftrightarrow2413$ for $m=4$ gives $3$ classes. The $4{,}755$ classes of $\Sym_8$ have sizes $2$ ($29$ classes), $4$ ($260$), $8$ ($4094$), $12$ ($10$), $16$ ($309$), $18$ ($1$), $20$ ($1$), $24$ ($39$), $32$ ($9$), $40$ ($1$), $48$ ($1$) and $56$ ($1$).
\end{proposition}

\begin{proof}
The program \texttt{known\_classes.py} applies union--find to the $m!$ permutations with the listed moves: for each permutation it generates its seven nontrivial symmetric images, and, when the permutation is a direct sum whose first summand is $\iota_k$, $\dd_k$ ($k\ge2$), $231$ or $312$, the permutation with that summand replaced; the counts are read off. The program \texttt{provenance.py} recomputes the classes by breadth-first search from the smallest element of each class and records, for each class, a spanning tree of moves connecting all members to it (for $\Sym_8$: $34{,}598$ symmetry moves, $873$ moves of type \eqref{eq:BWX} and $94$ of type \eqref{eq:SW}); the two computations agree. The class sizes are read from the output. The values $16$, $91$, $595$ agree with the classifications of $\Sym_5$, $\Sym_6$, $\Sym_7$ recorded in \cite{Vatter2026}; for $\Sym_4$ the theorem-generated partition has four blocks and the true classification three, the difference being $1342\sim2413$.
\end{proof}

The $56$-element class of $12345678$ is $\mathcal E_8\cup\mathcal E_8^r$ of \cref{thm:monoclass}.

\subsection{The lower bound: exact avoidance counts by the cluster expansion}\label{sec:counts}
For the lower bound we compute $|\Av_n(\tau)|$ for $n\le14$ and one representative $\tau$ of each $\equiv_8$-class, using \cref{thm:cluster}: $|\Av_n(\tau)|=n!-\sum_{r=8}^{n}\binom nr^2(n-r)!\,c_r(\tau)$, and $c_r(\tau)$ is the sum of the cluster weights of all clusters of $\tau$ of size $r$. The program \texttt{clusters.cpp} proceeds in three steps.

\emph{Generation.} Clusters are generated level by level, starting from $\tau$ itself. Given a cluster $\rho'$ of size $r'$ and $1\le s\le a$, choose a set $P$ of $a-s$ positions of $\tau$ (the roles that the new occurrence will share with $\rho'$), an embedding of the pattern $\std(\tau|_P)$ into $\rho'$ (the shared points), and an interleaving of the $s$ remaining points of $\tau$ among the points of $\rho'$, in positions and in values, that is consistent with the new occurrence being $\tau$; the result is a permutation of size $r'+s$ every point of which lies in an occurrence, i.e. a cluster. Every cluster $\sigma$ of size at most $a+K$ is obtained: order a minimal set of occurrences covering $\sigma$ and take successive unions; each union is a cluster of size at most $|\sigma|$ obtained from the previous one by exactly one such step with $s\ge1$. Generated permutations are stored in a hash set keyed by an injective packing of the permutation, so each cluster is produced once.

\emph{Weights.} For each cluster $\sigma$ of size $r$, the indicator $T\mapsto\bb{\sigma|_T\supseteq\tau}$ on subsets of $[r]$ is computed by a bit-parallel superset transform from the occurrences of $\tau$ in $\sigma$, and $c(\sigma)$ is its signed sum \eqref{eq:cw}. Summing over the clusters of size $r$ gives $c_r(\tau)$ exactly, since non-clusters have weight zero (\cref{thm:cluster}).

\emph{Assembly.} The avoidance counts are assembled from \eqref{eq:clusterexp} in $128$-bit integer arithmetic.

All arithmetic is exact. The computation for one pattern of length $8$ takes about one second for $n\le12$, eight seconds for $n\le13$ and five minutes for $n=14$ on one core. The counts were computed for all $4{,}755$ representatives for $n\le12$; for the $1{,}035$ representatives in the $490$ groups not separated by $|\Av_n|$, $n\le12$, also for $n=13$; and for the $8$ representatives in the $4$ groups not separated for $n\le13$, also for $n=14$. The outcome is:

\begin{proposition}\label{prop:separation}
The adaptively truncated cluster vectors of the $4{,}755$
representatives separate every pair: each pair differs in a coordinate
computed for both. Consequently their full vectors
$(c_8(\tau),\dots,c_{14}(\tau))$ are pairwise distinct. The numbers of groups of representatives with equal $(c_8,\dots,c_n)$, equivalently with equal $|\Av_t|$ for $t\le n$, are $1,8,256,4210,4751,4755$ for $n=9,\dots,14$.
\end{proposition}

The certificate is the table of all $4{,}755$ rows (class, size, representative, $c_8,\dots,c_{14}$ where computed, $|\Av_8|,\dots,|\Av_{14}|$, and the least $n$ separating the class from all others), provided in the companion repository together with the programs that regenerate it.

\begin{proof}[Proof of \cref{thm:E-intro}]
By \cref{prop:classes} each Wilf class of $\Sym_8$ is a union of $\equiv_8$-classes, so there are at most $4{,}755$ Wilf classes. By \cref{prop:separation} the representatives of distinct $\equiv_8$-classes have distinct avoidance sequences, so distinct $\equiv_8$-classes lie in distinct Wilf classes. Hence the Wilf classes are exactly the $\equiv_8$-classes, and the equivalences generating $\equiv_8$ generate all Wilf equivalences of length $8$. The refinement numbers are those of \cref{prop:separation}; since two representatives agree at $n\le13$ and differ at $n=14$, the bound $14$ is least.
\end{proof}

\subsection{Cross-checks}\label{sec:crosschecks}
The correctness of \cref{prop:separation} rests on the program \texttt{clusters.cpp}. Besides the coverage argument above, the following checks were made; their logs are in the companion repository.
\begin{enumerate}
\item For $\tau=123$ the program returns the Catalan numbers, and for $\tau\in\{1234,1324,1342\}$ the known sequences, for $n\le9$.
\item For $\tau=12345678$ it returns $|\Av_n|=476591309$, $6162155981$, $85494566892$ for $n=12,13,14$, the values obtained from the hook-length formula (Schensted's theorem), and its cluster numbers $c_9,\dots,c_{14}=-16,153,-1140,7315,-42504,230230$ are the values predicted by \cref{thm:bessel}.
\item An independently written counter (\texttt{autom.cpp}, an automaton over ``active'' partial permutations, developed for a different purpose) returns the same $|\Av_{10}|$ for twelve representatives, including the eight hardest ones, and the same $|\Av_{11}|$ for the eight hardest.
\item For $40$ classes, a second member of the class -- reached from the representative by a move of type \eqref{eq:BWX} in $29$ cases, of type \eqref{eq:SW} in one case, and by a symmetry in ten cases -- has the same $c_8,\dots,c_{12}$ as the representative. This tests the implementation and the cited theorems together.
\item An earlier computation in the same research programme produced the same partition of $\Sym_8$ by two different exact algorithms: for $n\le12$, counts of distinct length-$8$ subpatterns of each host, weighted by the symmetry orbits of the hosts; for $n=13,14$, unions of bitmaps over the orders of a tail block of each host. Its class membership file, its $20{,}063$ certificate entries (all $4{,}755$ representatives at $n=9,\dots,12$, the $1{,}035$ representatives of the groups unresolved at $n=12$ at $n=13$, and the $8$ representatives of the four last groups at $n=14$) and its $161{,}280$ per-pattern counts of all $40{,}320$ patterns at $n=9,\dots,12$ were compared with the present certificate: the classes, the representatives and every count agree. The two computations share no code.
\item The four pairs of \cref{tab:hardest} are separated only by $|\Av_{14}|$, and the two algorithms above come from the same research programme. A third method, as naive as possible and written independently of both (\texttt{bruteforce\_contain.cpp}), enumerates every placement of an occurrence of $\tau$ in $[n]\times[n]$ together with every arrangement of the remaining $n-8$ points and records the resulting permutations in a bitmap indexed by rank, so that each permutation containing $\tau$ is counted once ($6.5\cdot10^9$ placements for $n=14$). It returns the certificate's values of $|\Av_{11}|$ and $|\Av_{12}|$ for all $4{,}755$ representatives and of $|\Av_{13}|$ and $|\Av_{14}|$ for the eight representatives of \cref{tab:hardest}, exactly.
\end{enumerate}

\subsection{The hardest pairs}\label{sec:hardest}
Since $c_9=-16$ for every pattern of length $8$, $c_{10}=160-j$ determines $j$, and the pair $(j,c_{11})$ determines $g_3$, \cref{prop:separation} says that the Ray--West parameter takes all $8$ values $0,\dots,7$ on $\Sym_8$, that $(j,g_3)$ takes $256$ values, and that counts through codimension four distinguish $4{,}210$
groups of candidate classes. Of these groups, $3{,}720$ are singletons;
the remaining $490$ groups contain $1{,}035$ representatives. Four pairs of classes are separated only at $n=14$; \cref{tab:hardest} lists them.

\begin{table}[ht]
\centering\small
\begin{tabular}{@{}lllrr@{}}
\toprule
shared $c_8,\dots,c_{13}$ & representative & structure & $c_{14}$ & $|\Av_{14}|$\\
\midrule
$1,-16,156,-1232,8921,-63684$ & $13426758$ & $1\oplus231\oplus231\oplus1$ & $470058$ & $85472873492$\\
 & $13427568$ & $1\oplus231\oplus312\oplus1$ & $470054$ & $85472873496$\\
\addlinespace
$1,-16,157,-1250,9014,-62345$ & $13678254$ & & $435117$ & $85461367267$\\
 & $32718564$ & & $433097$ & $85461369287$\\
\addlinespace
$1,-16,154,-1166,7688,-46234$ & $14567823$ & $1\oplus(12345\ominus12)$ & $254167$ & $85485717761$\\
 & $14567832$ & $1\oplus(12345\ominus21)$ & $254552$ & $85485717376$\\
\addlinespace
$1,-16,157,-1248,8990,-63406$ & $23815764$ & & $485282$ & $85460332594$\\
 & $25876314$ & & $486330$ & $85460331546$\\
\bottomrule
\end{tabular}
\smallskip
\caption{The four pairs of Wilf classes of $\Sym_8$ with equal $|\Av_n|$ for all $n\le13$. Each pair shares $c_8,\dots,c_{13}$ (first column) and is separated by $c_{14}$. Class sizes are $4$ for the first pair and $8$ for the others.}\label{tab:hardest}
\end{table}

Two of the four pairs have the form $1\oplus\alpha$, $1\oplus\beta$ with $\alpha\sim\beta$ in $\Sym_7$: $231\oplus231\oplus1=(312\oplus312\oplus1)^{i}$ and $312\oplus312\oplus1\sim231\oplus312\oplus1$ by \eqref{eq:SW}; and $12345\ominus12=(\dd_2\oplus\dd_5)^{r}$, $12345\ominus21=(\iota_2\oplus\dd_5)^{r}$ are equivalent by \eqref{eq:BWX}. Prepending a point destroys these equivalences -- the moves \eqref{eq:BWX} and \eqref{eq:SW} act on direct-sum prefixes only -- and the resulting patterns of length $8$ share the cluster numbers $c_9,\dots,c_{13}$ before separating at $c_{14}$ (by $4$ and by $385$ respectively). The cluster numbers are signed aggregates, and their agreement does not mean that the clusters themselves agree: for the first pair the numbers of clusters of size $11$ (permutations every point of which lies in an occurrence) are $1{,}942$ and $1{,}943$, while $c_{11}=-1232$ for both.

\subsection{Cutoffs for lengths four to eight}\label{sec:cutoff}
Let $W_m(n)$ be the number of classes of $\Sym_m$ under ``equal $|\Av_t|$ for all $t\le n$'', and let $\kappa(m)$ be the least $n$ at which $W_m(n)$ reaches the number of Wilf classes of $\Sym_m$. By \cref{thm:cluster}, $W_m(n)$ is the number of distinct vectors $(c_{m+1},\dots,c_n)$ on $\Sym_m$. \Cref{tab:cutoff} gives the values computed with the same programs, one representative per $\equiv_m$-class (for $m\le5$ the counts were also obtained with \texttt{autom.cpp}).

\begin{table}[ht]
\centering\small
\begin{tabular}{@{}lrrrrrrrr@{}}
\toprule
$m$ & $n=m{+}1$ & $m{+}2$ & $m{+}3$ & $m{+}4$ & $m{+}5$ & $m{+}6$ & Wilf classes & $\kappa(m)$\\
\midrule
$4$ & $1$ & $2$ & $3$ & $3$ & --- & --- & $3$ & $7$\\
$5$ & $1$ & $5$ & $14$ & $16$ & $16$ & --- & $16$ & $9$\\
$6$ & $1$ & $6$ & $58$ & $91$ & $91$ & --- & $91$ & $10$\\
$7$ & $1$ & $7$ & $154$ & $574$ & $595$ & --- & $595$ & $12$\\
$8$ & $1$ & $8$ & $256$ & $4210$ & $4751$ & $4755$ & $4755$ & $14$\\
\bottomrule
\end{tabular}
\caption{Refinement numbers $W_m(n)$ and cutoffs. The column $n=m+2$ is the number of distinct values of $j$ on $\Sym_m$. Dashes mark omitted entries, not zeros.}\label{tab:cutoff}
\end{table}

The family $P_m:=1\oplus(\iota_{m-3}\ominus12)$, $Q_m:=1\oplus(\iota_{m-3}\ominus21)$ realizes $2m-2$: for $m=4,\dots,8$ the two patterns have equal cluster numbers up to size $2m-3$ and different ones at size $2m-2$ ($c_6(1423)=46\ne45=c_6(1432)$; $c_8(14523)=-388\ne-389$; $c_{10}(145623)=3300\ne3301$; $c_{12}(1456723)=-29510\ne-29476$; $c_{14}(14567823)=254167\ne254552$). We have not found a proof that this persists for all $m$.

\section{Higher codimension: the inflation complex}\label{sec:higher}

Let $\rho\in\Sym_{a+k}$. The \emph{occurrence graph} $\Occ_\tau(\rho)$ has as vertices the occurrences of $\tau$ in $\rho$, two being adjacent when they differ in exactly one point (a \emph{one-move}); let $\bz^{(k)}(\rho)$ be its number of components, the number of one-move classes. For $k=2$ this is the line graph of the deletion graph and $\bz^{(2)}(\rho)=\bz(\Gam_\tau(\rho))$. Recall from \cref{sec:universal} the inflation counts $I_\tau(\sigma)$ and the quantity $\chi_\tau(\rho)=\sum_T(-1)^{|T|-a}I_\tau(\rho|_T)$ of \eqref{eq:chi}, whose sum over $\rho$ is the universal constant $C_k(a)$.

\subsection{The complex}
An \emph{inflation structure} in $\rho$ is a set $T\subseteq[a+k]$ together with the partition $T=B_1\sqcup\dots\sqcup B_a$ into the blocks of a decomposition $\rho|_T=\tau[\pi_1,\dots,\pi_a]$, the block $B_i$ playing the role of the $i$th point of $\tau$. Choosing one point in each block gives an occurrence of $\tau$, and these $\prod|B_i|$ occurrences are the vertices of a cell $\Delta(B_1)\times\dots\times\Delta(B_a)$, a product of simplices of dimension $|T|-a$. Restricting each block to a nonempty subset gives again an inflation structure (the restricted blocks remain intervals of the restricted pattern, with the same quotient $\tau$), so the faces of a cell are cells, and a cell is determined by its vertex set, the block $B_i$ being the set of $i$th points of its vertices. 

A point of $\rho$ may carry different roles in different occurrences, but a role is attached to a pair (point, occurrence): the role of a point in an occurrence $O$ is its index in $O$. Hence if an occurrence $O$ is a vertex of two cells with blocks $(B_i)$ and $(B'_i)$, its $i$th point lies in $B_i\cap B'_i$ for every $i$, and the common vertices of the two cells are exactly the selections from $(B_i\cap B'_i)$, the vertex set of a common face. Gluing the cells along common faces therefore produces a regular cell complex $K_\tau(\rho)$, the \emph{inflation complex}, whose $d$-cells are counted by $\sum_{|T|=a+d}I_\tau(\rho|_T)$; thus $\chi_\tau(\rho)$ is its Euler characteristic.

The vertices of $K_\tau(\rho)$ are the occurrences of $\tau$, and its edges are the inflation structures with a single block $\{x,x'\}$ of size two: pairs of occurrences $O,O'$ differing in one point such that $x$ and $x'$ form a bond of $\rho|_{O\cup O'}$ and carry the same role. Not every one-move is an edge. For $\tau=12$ and $\rho=123$ the three occurrences $\{1,2\},\{1,3\},\{2,3\}$ form a triangle in $\Occ_\tau(\rho)$, but $K_{12}(123)$ has only the two edges $12[1,12]$ and $12[12,1]$, joining $\{1,2\}$ to $\{1,3\}$ and $\{1,3\}$ to $\{2,3\}$; the one-move from $\{1,2\}$ to $\{2,3\}$ changes the role of the shared point $2$ and is not an edge. The two graphs nevertheless have the same components.

\begin{lemma}\label{lem:skeleton}
The components of the $1$-skeleton of $K_\tau(\rho)$ are exactly the one-move classes of occurrences of $\tau$ in $\rho$. Consequently $\chi_\tau(\rho)=\bz^{(k)}(\rho)$ whenever every component of $K_\tau(\rho)$ has Euler characteristic one.
\end{lemma}

\begin{proof}
Every edge is a one-move. Conversely let $O=Q\cup\{x\}$ and $O'=Q\cup\{x'\}$ be occurrences differing in one point, and let $\sigma:=\rho|_{O\cup O'}\in\Sym_{a+1}$. Deleting $x'$ or $x$ from $\sigma$ leaves an occurrence of $\tau$, so $x,x'\in\Del(\sigma)$, which by \cref{lem:run} is a run $z_1,\dots,z_m$ of $\sigma$. For consecutive points $z_i,z_{i+1}$ of the run, $\{z_i,z_{i+1}\}$ is a bond of $\sigma$ whose contraction $\sigma\setminus z_i$ is $\tau$; so $O\cup O'$ with the block $\{z_i,z_{i+1}\}$ and singleton blocks is an inflation structure, i.e.\ an edge of $K_\tau(\rho)$ joining the occurrences $(O\cup O')\setminus z_{i+1}$ and $(O\cup O')\setminus z_i$. Walking along the run from $x'$ to $x$ joins $O=(O\cup O')\setminus x'$ to $O'=(O\cup O')\setminus x$ by edges. The last assertion holds because the Euler characteristic of a finite complex is the sum over its components.
\end{proof}

\subsection{Cluster weights and inflation weights}
The universal constants isolate the inflation structures inside $c_r(\tau)$; the following identity says exactly what the remainder is.

\begin{proposition}\label{prop:weights}
For every $\sigma\in\Sym_r$ containing $\tau$,
\begin{equation}\label{eq:cvsI}
c(\sigma)=(-1)^{r-a}I_\tau(\sigma)+\sum_{T\subseteq[r]:\ \sigma|_T\supseteq\tau}(-1)^{r-|T|}\bigl(1-\chi_\tau(\sigma|_T)\bigr),
\end{equation}
while $c(\sigma)=I_\tau(\sigma)=0$ if $\sigma$ avoids $\tau$. Consequently $c_{a+d}(\tau)=(-1)^dN_d(a)+D_{a+d}(\tau)$ with
\begin{equation}\label{eq:D}
D_{a+d}(\tau):=\sum_{\sigma\in\Sym_{a+d}}\ \sum_{T:\ \sigma|_T\supseteq\tau}(-1)^{a+d-|T|}\bigl(1-\chi_\tau(\sigma|_T)\bigr),
\end{equation}
and $D_{a+2}(\tau)=-j(\tau)$.
\end{proposition}

\begin{proof}
Fix $\sigma\in\Sym_r$. For $S\subseteq[r]$, \eqref{eq:mobius} gives $\bb{\sigma|_S\supseteq\tau}=\sum_{T\subseteq S}c(\sigma|_T)$, and \eqref{eq:chi} applied to $\sigma|_S$ gives $\chi_\tau(\sigma|_S)=\sum_{T\subseteq S}(-1)^{|T|-a}I_\tau(\sigma|_T)$. Möbius inversion over the Boolean lattice of $[r]$ turns these into $c(\sigma)=\sum_T(-1)^{r-|T|}\bb{\sigma|_T\supseteq\tau}$ (the definition \eqref{eq:cw}) and $(-1)^{r-a}I_\tau(\sigma)=\sum_T(-1)^{r-|T|}\chi_\tau(\sigma|_T)$. Subtracting, and using $\chi_\tau(\sigma|_T)=0$ when $\sigma|_T$ avoids $\tau$ (the complex is empty), gives \eqref{eq:cvsI}. Summing over $\sigma\in\Sym_{a+d}$ and using $\sum_\sigma I_\tau(\sigma)=N_d(a)$ (proof of \cref{thm:universal}) gives the expansion of $c_{a+d}$, and $D_{a+2}=-j$ is \cref{prop:firstclusters} with $N_2(a)=2a(a+2)$.
\end{proof}

In codimension two the correction can be read off host by host.

\begin{proposition}\label{prop:chi2}
For $\rho\in\Sym_{a+2}$ every component of $K_\tau(\rho)$ has Euler characteristic one, so $\chi_\tau(\rho)=\bz(\Gam_\tau(\rho))$; and for $\rho$ containing $\tau$,
\begin{equation}\label{eq:c2}
c(\rho)=I_\tau(\rho)-\bigl(\bz(\Gam_\tau(\rho))-1\bigr).
\end{equation}
\end{proposition}

\begin{proof}
Write $\Gam=\Gam_\tau(\rho)$, $E=E(\Gam)$ and $V=V(\Gam)$. The vertices of $K=K_\tau(\rho)$ are the edges of $\Gam$, so $f_0=|E|$. A $1$-cell is an inflation structure on $T=[a+2]\setminus\{z\}$ with one block $\{y,y'\}$ of size two: then $\rho\setminus\{z,y\}=\rho\setminus\{z,y'\}=\tau$, so $y,y'\in R_z$, and $\{y,y'\}$ is a bond of $\rho\setminus z$, so $y,y'$ are consecutive in the run $R_z$ (\cref{lem:Rx}); conversely every consecutive pair of $R_z$ is such a block. Hence the $1$-cells with $T=[a+2]\setminus\{z\}$ number $\deg_\Gam(z)-1$ for $z\in V$, and $f_1=2|E|-|V|$. The $2$-cells are the inflation structures on $[a+2]$, which are exactly the elements of $T(\rho)\cup S(\rho)$ (\cref{def:cycles} and the discussion following it), so $f_2=t(\rho)+s(\rho)=\bone(\Gam)$ by \cref{lem:basis}. Therefore $\chi(K)=|E|-(2|E|-|V|)+\bone(\Gam)=|V|-|E|+\bone(\Gam)=\bz(\Gam)$.


\emph{Componentwise count.} The same count applies to each component. By \cref{lem:skeleton} the components of $K$ are the one-move classes, i.e.\ the edge sets of the components of $\Gam$ restricted to $V$; the $1$-cells at $z$ and the $2$-cells through $z$ belong to the component of $z$; and the elements of $T(\rho)\cup S(\rho)$ lying in a component $\Gam_c$ form a basis of the cycle space of $\Gam_c$, because the cycle space of $\Gam$ is the direct sum of those of its components and each basis element of \cref{lem:basis} lies in one component. Hence $\chi(K_c)=|V_c|-|E_c|+\bone(\Gam_c)=1$ for every component. Finally \eqref{eq:c2} is \eqref{eq:cvsI} for $r=a+2$: the sub-hosts $\rho|_T$ with $|T|=a$ have a single occurrence, those with $|T|=a+1$ have a path of occurrences (\cref{lem:run}), so $\chi_\tau(\rho|_T)=1$ whenever $\rho|_T\supseteq\tau$ and $|T|\le a+1$, and the only surviving term is $T=[a+2]$.
\end{proof}

Summing \eqref{eq:c2} over $\rho$ recovers $c_{a+2}=2a(a+2)-j$ from \cref{thm:A}. The identity also shows that the correction $D_{a+d}$ is \emph{not} supported on the hosts that fail to be inflations of $\tau$: it lives on the hosts having a sub-host whose complex has Euler characteristic different from one, and such a host may well be an inflation. For $\tau=132$ the host $\rho=13254=132[132,1,1]$ has $I_\tau(\rho)=1$, but its occurrence $\{1,2,3\}$ and its three occurrences containing $\{4,5\}$ lie in different one-move classes, so $c(\rho)=1-(2-1)=0$; for $\tau=12$ the host $\sigma=14523=12[1,3412]$ has $I_{12}(\sigma)=1$ and $c(\sigma)=0$, the deviation coming from the sub-host $3412$, which contains $12$ in two disjoint occurrences. (A direct count for $\sigma=14523$: the position sets of sizes $2,3,4,5$ containing $12$ number $6,10,5,1$, and $-6+10-5+1=0$.) Summed over $\Sym_5$, the inflation hosts of $12$ contribute $-70$ to $c_5(12)=-56$ against $-N_3(2)=-72$, and the other hosts contribute $14$.

\subsection{Codimensions three and four, and the boundary at five}
The number of one-move classes is not universal in general (\cref{prop:k5} below), but there is a natural class of patterns for which the complex is always contractible and $g_k(\tau)=C_k(a)$ exactly. For $\tau\in\Sym_a$ put
\begin{equation}\label{eq:dinf}
d_\infty(\tau):=\min_{i<j}\max\{\,j-i,\ |\tau_j-\tau_i|\,\},
\end{equation}
the smallest Chebyshev distance between two points of $\tau$ (with $d_\infty=\infty$ for $a\le1$), and let $H_\tau$ be the multigraph on $[a]$ with red edges $\{i,i+1\}$ and blue edges $\{\tau^{-1}(v),\tau^{-1}(v+1)\}$; a red and a blue edge with the same ends form a cycle of length two, and a forest has infinite girth.

\begin{lemma}[projection-preserving elimination]\label{lem:elimination}
Let $F$ be a nonempty finite set of assignments to variables whose
candidate values are planar points with injective coordinates. Suppose
feasibility is specified by unary restrictions and order constraints,
each requiring a list of distinct variables to be strictly increasing in
one coordinate. Assume the incidence graph between
variables and nonunary constraints is a forest.
The complex whose cells are products $\prod_i\Delta(S_i)$ with every
selection from the $S_i$ in $F$ is contractible.
\end{lemma}

\begin{proof}
At each stage some variables have been fixed. Let $F'$ be the set of
remaining feasible assignments, and define each current candidate set
$D_i$ to be its projection onto variable $i$. We maintain the following
invariant when fixing $i$ to $d\in D_i$:
\begin{equation}\label{eq:projection-invariant}
\operatorname{pr}_{-i}F'
 =\operatorname{pr}_{-i}\{x\in F':x_i=d\}.
\end{equation}
Thus no feasible assignment of the other variables is lost.
Bounds from fixed variables and constraints with only one remaining
variable are unary restrictions and are absorbed into the candidate
sets. Splitting an ordered list at a fixed variable leaves shorter
ordered lists and preserves the forest property.

In a component containing a nonunary constraint, choose a pendant
constraint: all but at most one of its variables occur in no other
nonunary constraint. Such a constraint exists in a finite incidence
forest, by deleting private leaf variables and then taking a leaf
constraint of the remaining tree. Call the possible shared variable the
articulation. In the order of this constraint, fix the private variables
before the articulation from left to right, each to its least current
candidate. Fix those after it from right to left, each to its greatest
current candidate. If there is no articulation, use successive least
candidates throughout.

For a least-candidate step, replacement can only decrease the selected
coordinate, so it preserves the inequality with the next variable.
The preceding variable is already fixed; the replacement is above it
because it belongs to the projection of $F'$ and hence occurs in some
feasible assignment with that same predecessor. The variable is private,
so no other nonunary constraint changes, and its unary restrictions hold
by the definition of $D_i$. This proves \eqref{eq:projection-invariant};
the greatest-candidate step is symmetric. In particular all feasible
articulation values are preserved. A variable in no nonunary constraint
can be fixed to any current candidate.

Each step is also a strong deformation retraction of the complex.
If $d$ can replace variable $i$ in every feasible assignment, then
replacing $S_i$ by $S_i\cup\{d\}$ preserves the all-selections
condition. In this enlarged cell, move the $i$th barycentric coordinate
linearly to the vertex $d$, leaving the other coordinates unchanged.
These maps agree on common faces and fix the subcomplex where $i=d$.
Repeating the elimination removes at least one variable at each step,
retains nonemptiness by \eqref{eq:projection-invariant}, and ends at a
single assignment.
\end{proof}

\begin{theorem}[separated patterns]\label{thm:separated}
If $d_\infty(\tau)>k$ and $H_\tau$ has girth greater than $k$, then $K_\tau(\rho)$ is contractible for every $\rho\in\Sym_{a+k}$ containing $\tau$, and consequently $g_k(\tau)=C_k(a)$. In particular $g_3(\tau)=C_3(a)$ whenever $d_\infty(\tau)\ge4$, and $g_2(\tau)=C_2(a)$, i.e. $j(\tau)=0$, whenever $d_\infty(\tau)\ge3$.
\end{theorem}

\begin{proof}
\emph{Unique roles.} Label the points of an occurrence of $\tau$ in $\rho$ by their roles $1,\dots,a$. A point of $\rho$ at position $p$ and value $v$ carrying the role $i$ in some occurrence satisfies $p-k\le i\le p$ and $v-k\le\tau_i\le v$, since an occurrence omits only $k$ points. If a point carried two roles $i\ne i'$ in different occurrences, then $|i-i'|\le k$ and $|\tau_i-\tau_{i'}|\le k$, contradicting $d_\infty(\tau)>k$. So every point of $\rho$ used by some occurrence has a unique role; let $D_i$ be the set of points carrying the role $i$. The $D_i$ are disjoint, and if $U:=\{i:|D_i|\ge2\}$ then $a+|U|\le\sum_i|D_i|\le a+k$, so $|U|\le k$; the roles outside $U$ are carried by a fixed point in every occurrence.

\emph{A forest of order constraints.} An assignment of a point of $D_i$ to each role $i$ is an occurrence if and only if the chosen points are in the right relative order. Partition $U$ into maximal runs of consecutive roles and into maximal runs of roles with consecutive values of $\tau$. Two roles in $U$ lying in different position runs are separated by a fixed role, and every candidate point occurs in some occurrence together with that fixed point, so every pair of candidates for the two roles is already in the right horizontal order; likewise for value runs. Hence the feasible assignments are exactly the choices from the $D_i$ that are increasing along each position run (in position) and along each value run (in value). 

Let $B$ be the bipartite multigraph whose vertices are the position runs and the value runs of $U$ and whose edges are the roles in $U$, each joining its two runs. If $U$ has $u$ roles, $r$ position runs, $s$ value runs and $H_\tau[U]$ has $c$ components, then $H_\tau[U]$ has $2u-r-s$ edges and $B$ has $u$ edges on $r+s$ vertices, and the components correspond, so both have cycle rank $u-r-s+c$; since $H_\tau[U]$ is a subgraph of $H_\tau$ on at most $k$ vertices and $H_\tau$ has girth $>k$, $H_\tau[U]$ is a forest, and so is $B$.

\emph{Contraction.} The cells of $K_\tau(\rho)$ are precisely the
products $\prod_i\Delta(S_i)$, with $S_i\subseteq D_i$ nonempty,
for which every selection is feasible. Indeed, the horizontal and
vertical comparisons for all selections say that the blocks $S_i$
are separated in both coordinates, hence give an inflation structure.
The variable--constraint incidence graph is obtained by subdividing
the edges of $B$, so it is a forest. Apply \cref{lem:elimination} to
these position and value order constraints. Its projection invariant
ensures that eliminating a private role preserves every feasible
assignment of all remaining roles, not just their individual domains.
The complex is therefore contractible.

Thus every nonempty $K_\tau(\rho)$ has Euler characteristic one, and \eqref{eq:gkC} gives $g_k(\tau)=C_k(a)$. For $k=3$ the girth condition is implied by $d_\infty\ge4$: a red-blue $2$-cycle joins two roles at Chebyshev distance one, and a triangle in the union of the two paths contains two adjacent edges of one colour, whose end roles are at distance two in that colour and at distance one in the other. For $k=2$ the girth condition is implied by $d_\infty\ge3$ in the same way.
\end{proof}

There is a downward counterpart to this separation principle.
For $1\le k<a$, all $\binom ak$ positional deletions of $k$ points
from $\tau$ give distinct patterns if and only if
$\min_{i<j}(|i-j|+|\tau_i-\tau_j|)\ge k+2$.
This criterion was stated by Homberger \cite[preprint, Thm.~19]{Homberger2012};
a corrected proof is given by Bevan, Homberger and Tenner
\cite[preprint, Thm.~2.25]{BHT2018}. Hegarty \cite{Hegarty2013} proved the
necessity direction. The criterion concerns distinct patterns below one
permutation and the taxicab metric. In contrast, \cref{thm:separated}
concerns occurrences of a fixed pattern in all larger hosts, uses the
Chebyshev metric together with a girth condition, and proves
contractibility of their inflation complexes.

Patterns with $d_\infty\ge3$ exist from length $9$ on ($369258147$) and patterns with $d_\infty\ge4$ from length $16$ on ($4,8,12,16,3,7,11,15,2,6,10,14,1,5,9,13$), by an exhaustive search for shorter lengths (companion repository); the pattern $\tau_i=1+(4(i-1)\bmod a)$ has $d_\infty(\tau)\ge4$ for every odd $a\ge17$, and $g_3=C_3(a)$ was confirmed for $a=17,19$ by the deletion-overlap formula of \cref{thm:boxes} and by direct enumeration. The hypothesis $d_\infty\ge4$ cannot be weakened to $d_\infty\ge3$ for $k=3$: the pattern $\tau_i=1+(3(i-1)\bmod 17)$ has $d_\infty=3$, $j=0$ and $g_3=C_3(17)-3$, and $\tau_i=1+(3(i-1)\bmod13)$ has $g_3=C_3(13)-3$ (by \cref{thm:boxes} and by two independent insertion counters, companion repository).

\begin{corollary}[cubic pattern dependence in codimension three]\label{cor:cubic-spread}
For every odd $a\ge17$, let
$\sigma^{(a)}_i=1+(4(i-1)\bmod a)$, $1\le i\le a$. Then
\[
g_3(\sigma^{(a)})-g_3(\iota_a)
  =(a+1)(a^2+2a-5).
\]
In particular, no polynomial in $a$ independent of the pattern can
agree with both families up to remainders of degree less than three.
\end{corollary}

\begin{proof}
Since $a$ is odd, multiplication by $4$ modulo $a$ is a bijection.
If two positions differ by $d\in\{1,2,3\}$, their values differ by
$4d$ or $a-4d$, both at least $4$ for $a\ge17$. Positions at distance
at least $4$ already have Chebyshev distance at least $4$.
Thus $d_\infty(\sigma^{(a)})\ge4$, and \cref{thm:separated} gives
$g_3(\sigma^{(a)})=C_3(a)$. Subtract the polynomial $M_3(a)$ in
\cref{thm:layered} to obtain the identity. A difference of two
polynomials of degree less than three cannot equal this cubic on
infinitely many odd integers.
\end{proof}

Ray and West \cite[p.~57]{RayWest2003} conjecture that the
pattern-dependent part of the fixed-codimension polynomial has degree
less than the codimension. Read as asserting common coefficients
in degrees at least $k$, that clause is false already for $k=3$ by
\cref{cor:cubic-spread}. The comparison does not challenge their
leading-term estimate \cite[Cor.~9.5]{RayWest2003}, nor does it settle the separate
question of polynomiality for other specified families of patterns.

\begin{theorem}[one-move classes in codimensions three and four; computational]\label{thm:k34}
For $k=3$ and every $\tau$ of length $a\le6$, and for $k=4$ and every $\tau$ of length $a\le5$, every component of $K_\tau(\rho)$ has Euler characteristic one for every $\rho\in\Sym_{a+k}$. Consequently $\chi_\tau(\rho)=\bz^{(k)}(\rho)$ for every $\rho$, $\sum_{\rho\in\Sym_{a+k}}\bz^{(k)}(\rho)=C_k(a)$, and, in these ranges,
\begin{equation}\label{eq:g3classes}
g_3(\tau)=C_3(a)-\Delta_3(\tau),\qquad \Delta_3(\tau):=\sum_{\rho\in\Sym_{a+3}}\bigl(\bz^{(3)}(\rho)-1\bigr),
\end{equation}
the sum over the $\rho$ containing $\tau$, and likewise for $g_4$ with $k=4$.
\end{theorem}

\begin{proof}
For $a=1$, $K_\tau(\rho)$ is the simplex on all points of $\rho$,
so the assertion holds directly. For $a\ge2$ in the stated ranges,
the program \texttt{chi\_components.cpp} enumerates all $\rho\in\Sym_{a+k}$ containing $\tau$, lists the occurrences of $\tau$ in $\rho$ as position sets, forms the one-move classes by union--find on the relation ``differ in one point'', and then enumerates every cell of $K_\tau(\rho)$: every $T\subseteq[a+k]$ with $|T|\ge a$ and every decomposition of $\rho|_T$ into $a$ consecutive blocks that are value intervals with quotient $\tau$. Each cell is assigned to the class of the occurrence formed by the first points of its blocks (a vertex of the cell, hence in the class of all its vertices by \cref{lem:skeleton}), and $(-1)^{|T|-a}$ is added to the Euler characteristic of that class. 

For all $872$ patterns of lengths $2$ through $6$ with $k=3$ and all $152$ patterns of lengths $2$ through $5$ with $k=4$, no class has Euler characteristic different from one; the class totals equal $C_k(a)$ (for instance $18{,}344$ classes and $42{,}336+31{,}104+7{,}776+664$ cells summed over the $\rho\in\Sym_9$ containing a given $\tau\in\Sym_6$). The class totals were also obtained by the independent program \texttt{classes\_k.cpp}, which computes $\sum_\rho\bz^{(k)}(\rho)$ without cells. The logs are in the companion repository. Then \eqref{eq:g3classes} is \eqref{eq:gkC} with $\chi_\tau(\rho)=\bz^{(3)}(\rho)$.
\end{proof}

We conjecture that \cref{thm:k34} holds for all $a$; possible approaches include computing the Euler characteristic of each component of $K_\tau(\rho)$ by a cell count as in \cref{prop:chi2}, or proving that each component is contractible as in \cref{thm:separated}. The following shows that such a statement cannot hold in all codimensions.

\begin{proposition}[codimension five]\label{prop:k5}
Let $\tau=123$. For the four permutations
\[
\rho\in\{34127856,\ 34172856,\ 35127846,\ 35172846\}
\]
one has $\chi_\tau(\rho)=0$ and $\bz^{(5)}(\rho)=1$, and for every $\tau\in\Sym_3$ one has $\sum_{\rho\in\Sym_8}\bz^{(5)}(\rho)=44892$ while $C_5(3)=44888$. For $a=4$ the sum $\sum_{\rho\in\Sym_9}\bz^{(5)}(\rho)$ is not even universal: it equals $302001$ for $\tau=1234$ and $\tau=1324$ but $301992=C_5(4)$ for $\tau=2413$.
\end{proposition}

\begin{proof}
Direct computation (\texttt{chi\_k.py}, \texttt{classes\_k.cpp}); the four permutations of length $8$ are the only $\rho\in\Sym_8$ with $\chi_{123}(\rho)\ne\bz^{(5)}(\rho)$.
\end{proof}

Thus the universal object in every codimension is the Euler characteristic $\chi_\tau$; the number of one-move classes agrees with it in codimensions at most two (\cref{prop:chi2}) and, in the computed ranges, three and four (\cref{thm:k34}), but not in codimension five. In particular the conjectural form $g_3=C_3(a)-\Delta_3$ with $\Delta_3$ a count of one-move classes is proved only for $a\le6$; the exact forms of $g_3$ in \cref{sec:codim3} do not depend on it.

\section{The increasing pattern and the discrete Bessel kernel}\label{sec:monotone}

For $\tau=\iota_a$ the generating function \eqref{eq:egf} is a Plancherel probability. Let $f^\lambda$ be the number of standard Young tableaux of shape $\lambda$ and let $P_\theta$ be the Poissonized Plancherel measure, $P_\theta(\lambda)=e^{-\theta}\theta^{|\lambda|}(f^\lambda/|\lambda|!)^2$ for partitions $\lambda$ of all sizes. 

By Schensted's theorem \cite{Schensted1961}, the length of the longest increasing subsequence of $\pi\in\Sym_n$ is the first row of the Robinson--Schensted shape of $\pi$, so $|\Av_n(\iota_a)|=\sum_{\lambda\vdash n,\ \lambda_1\le a-1}(f^\lambda)^2$, where $f^\lambda$ is computed by the hook-length formula \cite{FRT1954}, and \eqref{eq:egf} reads
\begin{equation}\label{eq:plancherel}
\sum_{r\ge a}c_r(\iota_a)\frac{\theta^r}{r!^2}=P_\theta(\lambda_1\ge a).
\end{equation}
Under $P_\theta$ the random set $\{\lambda_i-i:i\ge1\}\subset\ZZ$ is a determinantal point process whose correlation kernel is the discrete Bessel kernel: this is Theorem~2 of Borodin, Okounkov and Olshanski \cite{BOO2000}, whose configuration $\mathcal D(\lambda)=\{\lambda_i-i\}$ and kernel (1.9) we use unchanged, and Theorem~1.3 with (1.10) of Johansson \cite{Johansson2001}; the kernel is
\begin{equation}\label{eq:kernel}
K(x,y)=\sqrt\theta\,\frac{J_xJ_{y+1}-J_{x+1}J_y}{x-y}=\sum_{s\ge1}J_{x+s}\,J_{y+s}\qquad(x,y\in\ZZ;\ J_m:=J_m(2\sqrt\theta)),
\end{equation}
the series form being \cite[Prop.~2.9]{BOO2000} (for $x=y$ the quotient is read as a limit), where $J_m$ is the Bessel function of the first kind. For $m\ge0$,
$J_m(2u)=\sum_{k\ge0}(-1)^k u^{m+2k}/(k!\,(m+k)!)$;
negative integer indices are defined by $J_{-m}=(-1)^mJ_m$.
Only nonnegative indices occur in the restriction below. In these coordinates the point $\lambda_1-1$ is the largest point of the configuration, so $\lambda_1\ge a$ if and only if the configuration meets $\ZZ_{\ge a-1}$, and the gap probability of a determinantal process is a Fredholm determinant (compare \cite[Thm.~1.3]{Johansson2001}, where the same determinant appears as $P_\theta(\lambda_1\le n)=\det(I-K)_{\ell^2(\{n,n+1,\dots\})}$):
\begin{equation}\label{eq:fredholm}
P_\theta(\lambda_1\ge a)=1-\det(I-K)_{\ell^2(\ZZ_{\ge a-1})}=\sum_{p\ge1}(-1)^{p+1}\sum_{a-1\le x_1<\dots<x_p}\det\bigl[K(x_i,x_j)\bigr]_{i,j=1}^p .
\end{equation}
The representation \eqref{eq:plancherel}--\eqref{eq:fredholm} of the distribution of the longest increasing subsequence is classical, as is the older determinantal form of the generating function of $|\Av_n(\iota_a)|$ due to Gessel \cite{Gessel1990}; what is new below is the identification with the cluster generating function \eqref{eq:egf} and the coefficient extraction, which turns the Fredholm expansion into exact statements about the cluster numbers.

\begin{proposition}[first contribution of each Fredholm term]\label{prop:rectangle}
Let $a,p\ge1$, let $R=(a+p-1)^p$ be the rectangular partition with
$p$ rows of length $a+p-1$, and put $L=p(a+p-1)$. The unsigned
$p$-point term of \eqref{eq:fredholm} satisfies
\[
F_{a,p}(\theta):=
\sum_{a-1\le x_1<\cdots<x_p}\det\bigl(K(x_i,x_j)\bigr)_{i,j=1}^{p}
=\frac{(f^R)^2}{(L!)^2}\theta^L+O(\theta^{L+1}).
\]
Consequently its signed contribution first occurs in degree $L$,
with coefficient $(-1)^{p+1}(f^R)^2/(L!)^2$.
\end{proposition}

\begin{proof}
For a partition $\lambda$ let
$N_a(\lambda)=\#\{i\ge1:\lambda_i-i\ge a-1\}$.
The determinantal correlation formula gives
$F_{a,p}(\theta)=\EE_\theta\binom{N_a(\lambda)}p$.
Since $\lambda_i-i$ decreases strictly with $i$,
$N_a(\lambda)\ge p$ if and only if $\lambda_p\ge a+p-1$,
which is equivalent to containment of the rectangle $R$.
There are no contributing partitions of size less than $L$;
the unique one of size $L$ is $R$, and $N_a(R)=p$.
The Poissonized Plancherel weight
$e^{-\theta}\theta^{|\lambda|}(f^\lambda/|\lambda|!)^2$
therefore gives the displayed leading term. The coefficient calculation
is legitimate as a power-series calculation, since there are only finitely
many partitions of each size. It also holds analytically near zero:
$N_a(\lambda)\le |\lambda|$, and the total weight at size $n$
is $e^{-\theta}\theta^n/n!$, so the factorial moments are locally
absolutely summable. The sign is that in \eqref{eq:fredholm}.
\end{proof}

\begin{theorem}\label{thm:bessel}
For all $a\ge1$ and $r\ge a$, $c_r(\iota_a)$ equals $r!^2$ times the coefficient of $\theta^r$ in \eqref{eq:fredholm}. In particular:
\begin{enumerate}
\item[(i)] $c_{a+k}(\iota_a)=(-1)^k\binom{2a+2k-2}{k}$ for $0\le k\le a+1$;
\item[(ii)] $c_{2a+2}(\iota_a)=(-1)^a\binom{4a+2}{a+2}-\Cat(a+1)^2$;
\item[(iii)] for every $r\ge a$ the one-point term ($p=1$) of \eqref{eq:fredholm} contributes exactly $(-1)^{r-a}\binom{2r-2}{r-a}$ to $c_r(\iota_a)$, and the summed $p$-point term with $p\ge2$ is
$O(\theta^{p(a+p-1)})$ as $\theta\to0$, for fixed $a,p$.
This exponent is sharp; \cref{prop:rectangle} gives the leading coefficient.
\end{enumerate}
\end{theorem}

\begin{proof}
The first statement is \eqref{eq:plancherel} and \eqref{eq:fredholm}. Write $u=\sqrt\theta$.

\emph{One-point term.} $\sum_{x\ge a-1}K(x,x)=\sum_{x\ge a-1}\sum_{s\ge1}J_{x+s}^2=\sum_{m\ge a}(m-a+1)J_m(2u)^2$. The classical product formula
\[
J_m(z)^2=\sum_{k\ge0}(-1)^k\frac{(2m+2k)!}{k!\,((m+k)!)^2\,(2m+k)!}\Bigl(\frac z2\Bigr)^{2m+2k}
\]
gives $[\theta^r]J_m(2\sqrt\theta)^2=(-1)^{r-m}\binom{2r}{r-m}/r!^2$. Hence the one-point term contributes to $c_r(\iota_a)$
\[
\sum_{m=a}^{r}(m-a+1)(-1)^{r-m}\binom{2r}{r-m}=\sum_{k=0}^{N}(N-k+1)(-1)^k\binom{2r}{k},\qquad N:=r-a .
\]
Using $\sum_{k=0}^{N}(-1)^k\binom nk=(-1)^N\binom{n-1}N$ and $\sum_{k=0}^{N}(N-k)(-1)^k\binom nk=\sum_{j=0}^{N-1}\sum_{k=0}^{j}(-1)^k\binom nk=(-1)^{N-1}\binom{n-2}{N-1}$ with $n=2r$, the sum equals $(-1)^N\bigl[\binom{2r-1}N-\binom{2r-2}{N-1}\bigr]=(-1)^{N}\binom{2r-2}{N}$. This is (iii) for $p=1$.

\emph{Higher terms.} The independent rectangle calculation in
\cref{prop:rectangle} gives the remaining assertion of (iii).
The two-point term starts in degree $2a+2$, and terms with $p\ge3$
start no earlier than degree $3a+6$. Thus only the one-point term
contributes in degrees at most $2a+1$, proving (i).
In degree $2a+2$, the two-point term contributes exactly
$-(f^{(a+1,a+1)})^2$ after multiplication by $(2a+2)!^2$.
The hook-length formula gives $f^{(a+1,a+1)}=\Cat(a+1)$,
which proves (ii).
\end{proof}

\begin{corollary}\label{cor:LIS}
For $a\le n\le2a+1$,
\[
\#\{\pi\in\Sym_n:\ \pi\text{ has an increasing subsequence of length }a\}=\sum_{r=a}^{n}(-1)^{r-a}\binom nr^2(n-r)!\binom{2r-2}{r-a}.
\]
\end{corollary}

\begin{proof}
Insert \cref{thm:bessel}(i) into \eqref{eq:clusterexp}.
\end{proof}

For $n=a+1$ the sum is $(a+1)^2-2a=a^2+1$; for $n=a+3$
it gives $g_3(\iota_a)$, as in \eqref{eq:M3}. For fixed $k=n-a$, the
right-hand side is a polynomial $P_k(a)$ of degree $2k$.
Ray and West \cite[Prop.~9.6, p.~87]{RayWest2003} had already
obtained polynomiality from Robinson--Schensted and the hook-length
formula. Ekhad, Shar and Zeilberger \cite[p.~2]{ESZ2015} already obtain such
polynomials from Robinson--Schensted and the hook-length formula,
explicitly state their validity for $a\ge k-1$, and compute them for
$k\le30$. In their notation this range is $d\ge r-1$, or $n\le2d+1$.
\Cref{cor:LIS} supplies one binomial-sum expression valid for arbitrary
$k$, rather than a list of polynomials computed separately.
\Cref{thm:bessel}(ii) gives the exact first defect:
\[
g_k(\iota_{k-2})=P_k(k-2)-\Cat(k-1)^2\qquad(k\ge3).
\]
Thus the already stated validity range is sharp. The formulas agree
with the polynomials displayed in \cite{ESZ2015} for $k\le5$.

\begin{corollary}[the next cluster coefficient]\label{cor:bessel-next}
For every $a\ge1$,
\[
c_{2a+3}(\iota_a)=(-1)^{a+1}\binom{4a+4}{a+3}
 +(a+1)\Cat(a+2)^2.
\]
\end{corollary}

\begin{proof}
Put $b=a+1$ and $n=2b+1$. By \cref{prop:rectangle}, terms with
$p\ge3$ begin in degree at least $3a+6>n$. The condition
$N_a(\lambda)\ge2$ is $\lambda_2\ge b$. The unique such partition
of size $2b$ is $R=(b,b)$, and the only such partitions of size
$n$ are $\mu=(b+1,b)$ and $\nu=(b,b,1)$; all three have
$N_a=2$. Thus
\[
F_{a,2}(\theta)=e^{-\theta}\left(
 \frac{\Cat(b)^2}{((2b)!)^2}\theta^{2b}
 +\frac{(f^\mu)^2+(f^\nu)^2}{(n!)^2}\theta^n
 +O(\theta^{n+1})\right).
\]
The hook-length formula gives
\[
f^\mu=\frac{2(2b+1)!}{(b+2)!\,b!}=\Cat(b+1),\qquad
f^\nu=\frac{b(b+1)(2b+1)!}{(b+2)!\,(b+1)!}
 =\frac b2\Cat(b+1).
\]
Since the sign of the two-point term is negative, its contribution
to $c_n(\iota_a)$ is
\[
n^2\Cat(b)^2-\left(1+\frac{b^2}{4}\right)\Cat(b+1)^2
=b\Cat(b+1)^2,
\]
where we used $\Cat(b+1)=2(2b+1)\Cat(b)/(b+2)$.
Adding the one-point contribution from \cref{thm:bessel}(iii)
proves the formula.
\end{proof}

\begin{remark}[verification]\label{rem:besseldata}
The script \texttt{bessel\_clusters.py} expands \eqref{eq:fredholm} exactly in rational arithmetic (terms with $p\le4$) and reproduces the cluster numbers $c_r(\iota_a)$ obtained by inverting \eqref{eq:clusterexp} with the hook-length values of $|\Av_n(\iota_a)|$, for $a\le5$ and $r\le2a+7$; the script \texttt{monotone\_clusters.py} checks the binomial law and its failure at $k=a+2$ for $a\le12$, $r\le a+16$. The next coefficient is proved in \cref{cor:bessel-next};
\texttt{check\_monotone\_coefficients.py} also checks it for $1\le a\le12$
using tableau counts obtained by corner removal rather than hook lengths.
The leading rectangular contributions are proved for all $p$ in
\cref{prop:rectangle}; for example, at $a=1$, $p=3$, their numerator is
$(f^{(3,3,3)})^2=1764$.
\end{remark}

\section{Further questions}\label{sec:questions}

\begin{enumerate}
\item \emph{A statistic for $g_3$.} Can the deletion-overlap formula be
replaced by a compact statistic of the intervals and diagram of $\tau$?
The layered correction gives one model; \cref{cor:cubic-spread} shows
that a general correction must allow cubic pattern dependence.
\item \emph{Components in codimensions three and four.} Does every
component of $K_\tau(\rho)$ have Euler characteristic one for $k=3,4$
and arbitrary pattern length? This would extend \cref{thm:k34} beyond
its finite range and interpret $C_k(a)-g_k(\tau)$ by extra one-move
classes. \Cref{prop:k5} rules out the same assertion for all $k$.
\item \emph{Distinguishing lengths.} For $m\ge4$, do the Wilf-class
cutoffs satisfy $2m-2\le\kappa(m)\le2m-1$? The lower-bound family
$P_m,Q_m$ in \cref{sec:cutoff} and the computed cutoffs through $m=8$
suggest these bounds; neither is established in general.
\end{enumerate}

\appendix
\section{Proof of the layered formula}\label[appendix]{app:layered}

Throughout, $\alpha=\dd_{b_1}\oplus\dots\oplus\dd_{b_t}$ is layered of length $a\ge2$ and $\beta:=\alpha\oplus1$. Layered permutations are involutions, so for a set $X$ of permutations closed under inversion we may freely pass between statements about positions and about values. For a permutation $\pi$ write $\pi\setminus\!\max$ and $\pi\setminus\!\mathrm{last}$ for the patterns obtained by deleting its maximum and its last point. We first prove the correction term in \eqref{eq:layered} for patterns ending in a singleton, by a recurrence in $a$, and then remove that restriction.

\subsection{Two states}
Let $p=q\oplus1$ be any pattern ending in its maximum and $N\ge2$. For $\pi\in\Sym_N$ with $\pi(N)\ne N$ call $\pi|_{[N-1]}$ and $\pi\setminus\!\max$ its two \emph{states}.

\begin{lemma}\label{lem:corner}
For $\pi\in\Sym_N$ with $\pi(N)\ne N$, $\pi$ avoids $p$ if and only if both states avoid $p$. Consequently
\begin{equation}\label{eq:cornerid}
|\Av_{N-1}(q)|=|\Av_N(p)|-\#\{\pi\in\Sym_N:\ \pi(N)\ne N,\ \text{both states avoid }p\}.
\end{equation}
\end{lemma}

\begin{proof}
An occurrence of $p$ in $\pi$ that avoids the last point lies in the first state; one that avoids the maximum lies in the second; and one that uses both is impossible, because its last point would be the last point of $\pi$, whose value is not $N$, whereas the last point of $p$ is its maximum. For \eqref{eq:cornerid}, the $\pi\in\Av_N(p)$ with $\pi(N)=N$ are exactly $\pi'\oplus1$ with $\pi'\in\Av_{N-1}(q)$.
\end{proof}

\begin{lemma}\label{lem:twostates}
Let $b_{N-1}(p):=\#\{\sigma\in\Sym_{N-1}:\sigma\supseteq p\}$ and
\[
J_N(p):=\#\{\pi\in\Sym_N:\ \pi(N)\ne N,\ \text{both states contain }p\}.
\]
Then $\#\{\pi\in\Sym_N:\pi(N)\ne N,\ \text{both states avoid }p\}=(N-1)(N-1)!-2(N-1)\,b_{N-1}(p)+J_N(p)$. Moreover $\pi\mapsto(\pi|_{[N-1]},\pi\setminus\!\max)$ is a bijection from $\{\pi\in\Sym_N:\pi(N)\ne N\}$ onto the pairs $(\sigma,\tau)\in\Sym_{N-1}^2$ with $\sigma\setminus\!\max=\tau\setminus\!\mathrm{last}$, so that
\begin{equation}\label{eq:Jsum}
J_N(p)=\sum_{\rho\in\Sym_{N-2}}f_p(\rho)\,f'_p(\rho),\qquad
\begin{aligned}
f_p(\rho)&:=\#\{\sigma\supseteq p:\sigma\setminus\!\max=\rho\},\\
f'_p(\rho)&:=\#\{\tau\supseteq p:\tau\setminus\!\mathrm{last}=\rho\},
\end{aligned}
\end{equation}
with $\sigma,\tau\in\Sym_{N-1}$; and if $p$ is an involution then $f_p(\rho)=f'_p(\rho^{-1})$.
\end{lemma}

\begin{proof}
There are $(N-1)(N-1)!$ permutations with $\pi(N)\ne N$; the first state takes each value in $\Sym_{N-1}$ exactly $N-1$ times (the last value of $\pi$ is free in $[N-1]$), and so does the second (the position of the maximum is free in $[N-1]$); inclusion--exclusion gives the count. A pair $(\sigma,\tau)$ with $\sigma\setminus\!\max=\tau\setminus\!\mathrm{last}$ determines $\pi$: the position of $N$ in $\pi$ is the position of $N-1$ in $\sigma$, and $\pi(N)=\tau(N-1)$; conversely the two states of $\pi$ satisfy the compatibility. The last assertion follows by inverting.
\end{proof}

Applying \cref{lem:corner,lem:twostates} to $p=\beta$ (an involution) with $N=a+3$ and $N=a+4$, and using $|\Av_M(\cdot)|=M!-g_{M-|\cdot|}(\cdot)$, gives
\begin{align}
g_2(\beta)&=g_2(\alpha)+2(a+2)\,g_1(\beta)-J_{a+3}(\beta),\label{eq:rec2}\\
g_3(\beta)&=g_3(\alpha)+2(a+3)\,g_2(\beta)-J_{a+4}(\beta).\label{eq:rec3}
\end{align}

\subsection{The correlation \texorpdfstring{$J_{a+3}(\beta)$}{J(a+3)}}

\begin{lemma}\label{lem:C}
Let $p$ be layered of length $r\ge2$ and let $\varepsilon(p)=1$ if $p$ has at least two layers and its last layer has size at least two, $\varepsilon(p)=0$ otherwise. Then $J_{r+2}(p)=2r^2+2r+1+\varepsilon(p)$.
\end{lemma}

\begin{proof}
Let $s:=p\setminus\!\mathrm{last}=p\setminus\!\max$ (the two deletions agree because the last layer of $p$ is decreasing). Put $E:=\{\rho\in\Sym_r:\rho\supseteq s\}$, so $|E|=(r-1)^2+1$ by \cref{thm:pratt}, and let $A\subseteq\Sym_r$ be the set of $r$ permutations obtained from $s$ by inserting, at any of the $r$ positions, a point whose value is that of the last point of $p$ (the maximum if the last layer of $p$ is a singleton, and otherwise the value immediately below the values of the last layer of $s$). Then $p\in A\subseteq E$.

\emph{Claim: $f'_p(\rho)=\bb{\rho\in E}+\bb{\rho\in A}+(r-1)\bb{\rho=p}$ for $\rho\in\Sym_r$.} Appending a point of rank $y\in[r+1]$ to $\rho$ gives a permutation containing $p$ if and only if $\rho=p$ (all $r+1$ values of $y$) or the appended point is the last point of an occurrence of $p$ whose other points form an occurrence of $s$ in $\rho$ with $y$ in the value gap prescribed by $p$. By \cref{lem:run} the occurrences of $s$ in $\rho$ are $\rho\setminus x$ for $x$ in one maximal run $R=\{x_1<\dots<x_m\}$ of $\rho$ (so $\rho\in E$ means $R\ne\emptyset$). For each $x\in R$ the admissible $y$ form an interval $G_x$ of ranks: the prescribed gap lies between the values $u_x<w_x$ of the two points of $\rho\setminus x$ that play the roles bounding it (with $u_x=-\infty$ or $w_x=+\infty$ for a gap at the bottom or the top), and $G_x$ consists of the ranks $y$ with $u_x<y\le w_x$ (a new point of rank $y$ lies above the old value $u$ if and only if $y>u$, and below the old value $w$ if and only if $y\le w$). Since the two bounding roles carry adjacent values in $s$, $G_x$ has one element, or two if the value of $x$ lies strictly between $u_x$ and $w_x$. The latter happens for some $x$ exactly when $\rho\in A$. 

The number of admissible $y$ is $|\bigcup_{x\in R}G_x|$, and we claim that this union has $1+\bb{\rho\in A}$ elements. The sets $G_x$ need not coincide (for $p=12$ and $\rho=21$ they are $\{2,3\}$ and $\{3\}$), but they are nested or adjacent: a role bounding the gap is played, in every occurrence $\rho\setminus x$, either by a fixed point of $\rho$ outside $R$, or by the point of $R\setminus x$ with a fixed index $i$ in the run, which is $x_{i+1}$ if $x\le x_i$ and $x_i$ otherwise, two points with consecutive values. If at most one endpoint of the gap varies with $x$, the intervals $G_x$ are nested and their union is the largest of them, of size $1+\bb{\rho\in A}$. 

If both endpoints vary, the two bounding roles are played by consecutive points of the run (their values are adjacent in $s$), so with the run increasing and the roles at indices $i,i+1$ the pairs $(u_x,w_x)$ are $(v_i,v_{i+1})$ for $x>x_{i+1}$, $(v_{i+1},v_{i+2})$ for $x<x_{i+1}$ and $(v_i,v_{i+2})$ for $x=x_{i+1}$, where $v_i<v_{i+1}<v_{i+2}$ are the consecutive values of $x_i,x_{i+1},x_{i+2}$; the union is $\{v_{i+1},v_{i+2}\}$, of size two, and $\rho\in A$ because $x_{i+1}$ lies in its gap (the decreasing case is symmetric). So the number of admissible $y$ is $\bb{\rho\in E}+\bb{\rho\in A}$ for $\rho\ne p$, and $r+1=1+1+(r-1)$ for $\rho=p$.

\emph{Claim: $|A\cap A^{*}|=2+\varepsilon(p)$}, where $X^*:=\{\rho^{-1}:\rho\in X\}$. Since $s$ is an involution, $A^*$ is the set obtained from $s$ by inserting a point at the prescribed position
$d=p_r$ with an arbitrary value: inversion interchanges the prescribed
value $d$ and the prescribed position $d$, so $\rho\in A\cap A^*$ means that deleting the point with the prescribed value and deleting the point at the prescribed position both give $s$. If the last layer of $p$ is a singleton, both deletions are the deletion of the maximum and of the last point, and by \cref{lem:run} they agree exactly when the maximum is last or the maximum and the last point form a decreasing run, i.e. $\rho=s\oplus1$ or $\rho$ enlarges the last decreasing layer of $s$: two permutations. 

If the last layer $\dd_{b_t}$ of $p$ has $b_t\ge2$, write $s=\gamma\oplus\dd_{b_t-1}$ with $|\gamma|=b:=r-b_t$; the prescribed value and position are both $b+1$. Either the point of value $b+1$ is at position $b+1$, or (\cref{lem:run}) it lies later and enlarges the final decreasing layer, or it lies earlier and enlarges the preceding decreasing layer, which is possible exactly when $\gamma$ is nonempty, i.e. $t\ge2$: $2+\varepsilon(p)$ permutations.
In each case deleting the point of value $b+1$ or the point at position
$b+1$ recovers $s$. Conversely, these are the three possible locations
of the prescribed value relative to the prescribed position; the run
condition fixes the insertion within the indicated layer. Thus the
listed constructions are distinct and exhaustive.

Now $E$ and $\{p\}$ are inversion-stable, $|A\cap E^*|=|A|=r$, and by \eqref{eq:Jsum}
\[
J_{r+2}(p)=\sum_\rho f'_p(\rho)f'_p(\rho^{-1})=|E|+2|A|+|A\cap A^*|+4(r-1)+(r-1)^2=2r^2+2r+1+\varepsilon(p).\qedhere
\]
\end{proof}

Since $\beta$ ends in a singleton, $J_{a+3}(\beta)=2(a+1)^2+2(a+1)+1$, and \eqref{eq:rec2} with $g_1(\beta)=(a+1)^2+1$ gives $g_2(\alpha\oplus1)=g_2(\alpha)+2(a+1)^3+1$, in agreement with $j=a-1$ for layered patterns (\cref{thm:boundary}).

\subsection{The correlation \texorpdfstring{$J_{a+4}(\beta)$}{J(a+4)}}
Put $n:=a+2$ and, for $\rho\in\Sym_n$ and $i=0,1,2$, let $A_i$ be the set of $\rho$ such that $\alpha$ occurs among the $a+i$ smallest values of $\rho$; let $E:=\{\rho\in\Sym_n:\rho\supseteq\beta\}$, $K:=E\cap A_0$ and $L:=|E\cap A_0\cap A_0^*|$. Then $A_0\subseteq A_1\subseteq A_2$, $E\subseteq A_1$ (an occurrence of $\beta$ minus its last point is an occurrence of $\alpha$ among the $a+1$ smallest values), $A_2$ and $E$ are inversion-stable, and $\varepsilon:=\varepsilon(\alpha)$.

\begin{lemma}\label{lem:f3}
For $\rho\in\Sym_n$, $f'_\beta(\rho)=\bb{\rho\in A_0}+\bb{\rho\in A_1}+\bb{\rho\in A_2}+(a+1)\bb{\rho\in E}-\bb{\rho\in K}$.
\end{lemma}

\begin{proof}
Append a point of rank $y\in[a+3]$ to $\rho$. If $\rho\in E$ every $y$ works, and the right-hand side is $\bb{A_0}+1+1+(a+1)-\bb{A_0}=a+3$. Otherwise the appended point must be the last point, hence the maximum, of the occurrence, so $\alpha$ must occur among the values below $y$; the admissible $y$ are those with $y-1\ge$ the least $m$ such that $\alpha$ occurs among the $m$ smallest values, and their number is $\bb{A_0}+\bb{A_1}+\bb{A_2}$.
\end{proof}

\begin{lemma}\label{lem:table}
$|A_0|=(a+2)(a+1)$, $|A_1|=(a+2)(a^2+1)$, $|A_2|=g_2(\alpha)=\tfrac12(a^4+2a^3+a^2+2a+6)$, $|E|=(a+1)^2+1$, $|K|=2a+3$, $|A_0\cap A_0^*|=7$, $|A_0\cap A_1^*|=4a+3$, $|A_1\cap A_1^*|=3a^2+2a+2+\varepsilon$, and $L=5$.
\end{lemma}

\begin{proof}
$|A_0|$: the $a$ smallest values form $\alpha$ and the two largest values are placed in any two of the $a+2$ positions in either order. $|A_1|$: the largest value is placed anywhere and the pattern of the other $a+1$ values is one of the $a^2+1$ permutations containing $\alpha$. $|A_2|=g_2(\alpha)$ by definition, and its value follows from $j(\alpha)=a-1$ (\cref{thm:boundary}). $|E|=g_1(\beta)$ by \cref{thm:pratt}.

\emph{$|K|$.}  $\rho\in A_0$ is $\alpha$ on the values $[a]$ with the two values $a+1,a+2$ inserted. An occurrence of $\beta$ using all $a$ small points needs one of the large points in the last position gap: $2(a+1)$ permutations ($2a$ with the other large point elsewhere, and two with both large points at the end in either order). If neither large point is in the last gap, an occurrence of $\beta$ must use both large points and omit one small point; its maximum $a+2$ must be its last point, so $a+2$ sits immediately before the last small point, which is the omitted one, and the point $a+1$ must sit where inserting a maximum into $\alpha$ minus its last point recovers $\alpha$; since $\alpha$ is layered this is the deletion of its maximum, and the position is unique. This gives one more permutation, so $|K|=2a+3$.

\emph{$|A_0\cap A_0^*|$.}  $\rho$ is $\alpha$ on the values $[a]$ and $\alpha$ on the positions $[a]$. Let $t\in\{0,1,2\}$ be the number of large values among the first $a$ positions. There are $\binom2t^2(2-t)!$ choices (which large values, which positions among the last two carry the small values displaced, and the order of the remaining large values), and the common $a-t$ points are consistent because the prefix and value restrictions of a layered pattern agree; the total is $2+4+1=7$.

\emph{$L$.} Among these seven, test containment of $\beta$: for $t=0$ both permutations contain $\beta$ (the last two values are large); for $t=1$, if the large value inside the prefix is $a+1$ then $a+2$ at the end completes the prefix occurrence in both arrangements, and if it is $a+2$ then $a+1$ completes the small-value occurrence exactly when it comes after the displaced small value: three permutations; for $t=2$ a $\beta$-occurrence omits one point, and omitting the value $a+2$ leaves $a+1$ too early while keeping it makes it the maximum in a non-final position: none. So $L=5$.

\emph{$|A_0\cap A_1^*|$.}  $\rho$ is $\alpha$ on the values $[a]$ and contains $\alpha$ in its first $a+1$ positions. If the last point is large: choose which large point is last and insert the other among the $a$ small points, $2(a+1)$ permutations. If the last point is small it is the last point of the small-value $\alpha$; let $s=\alpha\setminus\!\mathrm{last}=\alpha\setminus\!\max$; the prefix is $s$ with the two large points inserted and must contain $\alpha$. Let $c$ be the position gap of $s$ where inserting a new maximum recovers $\alpha$. A large point in gap $c$ gives $2a$ permutations. With neither large point in gap $c$, an occurrence of $\alpha$ must use both large points and omit one point of $s$. The run lemma then leaves exactly one configuration: if the last layer of $\alpha$ has size $\ge2$, both large points sit immediately after the first point of the shortened last layer, in decreasing order, and that point is omitted; if $\alpha$ ends in a singleton, the maximum sits immediately before the last point of $s$, which is omitted, and the other large point sits where $s$ is recovered from its own last-point deletion. In either case the indicated deletion recovers $s$, and the position of every inserted point is forced. This case therefore contributes one permutation. Altogether $|A_0\cap A_1^*|=2(a+1)+2a+1=4a+3$.

\emph{$|A_1\cap A_1^*|$.}  $\rho$ contains $\alpha$ after deleting its maximum and after deleting its last point. If the maximum is last the two deletions coincide and there are $g_1(\alpha)=a^2+1$ permutations. Otherwise the two deletions give a pair $(\sigma,\tau)$ of one-point extensions of $\alpha$ with $\sigma\setminus\!\max=\tau\setminus\!\mathrm{last}$, and by \cref{lem:twostates} each compatible pair arises from exactly one $\rho$ with non-final maximum; their number is $J_{a+2}(\alpha)=2a^2+2a+1+\varepsilon$ by \cref{lem:C}. Hence $|A_1\cap A_1^*|=3a^2+2a+2+\varepsilon$.
\end{proof}

\begin{lemma}\label{lem:J4}
$J_{a+4}(\beta)=\tfrac12\bigl(3a^4+22a^3+57a^2+66a+38\bigr)+\varepsilon$.
\end{lemma}

\begin{proof}
By \eqref{eq:Jsum} and \cref{lem:f3}, with $u:=\bb{A_0}+\bb{A_1}+\bb{A_2}$, $c:=a+1$, and $\langle g,h\rangle:=\sum_\rho g(\rho)h(\rho^{-1})$,
\[
J_{a+4}(\beta)=\langle u,u\rangle+2c\langle u,\bb E\rangle+c^2|E|-2\langle u,\bb K\rangle-2c\langle\bb E,\bb K\rangle+\langle\bb K,\bb K\rangle .
\]
Using $A_0\subseteq A_1\subseteq A_2=A_2^*$, $E=E^*\subseteq A_1$ and $K^*=E\cap A_0^*$: $\langle u,u\rangle=|A_0\cap A_0^*|+2|A_0\cap A_1^*|+|A_1\cap A_1^*|+2|A_0|+2|A_1|+|A_2|$; $\langle u,\bb E\rangle=|K|+2|E|$; $\langle u,\bb K\rangle=L+2|K|$; $\langle\bb E,\bb K\rangle=|K|$; $\langle\bb K,\bb K\rangle=L$. Hence
\[
J_{a+4}(\beta)=|A_0\cap A_0^*|+2|A_0\cap A_1^*|+|A_1\cap A_1^*|+2|A_0|+2|A_1|+|A_2|+(c^2+4c)|E|-4|K|-L,
\]
and substituting \cref{lem:table} gives the polynomial.
\end{proof}

\subsection{Conclusion}
Insert \cref{lem:J4} and $g_2(\beta)=\tfrac12\bigl((a+1)^4+2(a+1)^3+(a+1)^2+2(a+1)+6\bigr)$ into \eqref{eq:rec3}:
\begin{equation}\label{eq:rec3final}
g_3(\alpha\oplus1)-g_3(\alpha)=\tfrac12\bigl(2a^5+15a^4+40a^3+49a^2+42a+34\bigr)-\varepsilon(\alpha)=M_3(a+1)-M_3(a)-\varepsilon(\alpha).
\end{equation}
Appending a singleton to $\alpha$ makes its last layer an interior layer, so the number $I(\cdot)$ of interior layers of size $\ge2$ satisfies $I(\alpha\oplus1)=I(\alpha)+\varepsilon(\alpha)$. Both layered patterns of length $2$ have $g_3=5!-1=119=M_3(2)$ and $I=0$. Induction on $a$ with \eqref{eq:rec3final} proves $g_3(\tau)=M_3(|\tau|)-I(\tau)$ for every layered $\tau$ ending in a singleton, provided it holds for all layered patterns of smaller length; and a layered $\tau=\tau'\oplus\dd_b$ with $b\ge2$ is Wilf-equivalent to $\tau'\oplus\iota_b$ by \eqref{eq:BWX} applied after reverse-complement, which ends in a singleton, is layered, and has the same $I$ (the replaced layer was final and the new layers are singletons). This proves \eqref{eq:layered} for layered $\tau$, and reversal gives the reverse-layered case. \qed

\begin{remark}
The identities of \cref{lem:C,lem:f3,lem:table,lem:J4} and the recurrence \eqref{eq:rec3final} were also verified by exhaustive computation for every layered $\alpha$ of length at most $5$, and \eqref{eq:layered} for every layered pattern of length at most $8$ by the cluster program and by the deletion-overlap formula (companion repository).
\end{remark}

\section*{Data and code availability}
% BEGIN RELEASE AVAILABILITY
The code, certificates, row-level recounts and buildable source are
provided in the versioned companion at
\href{https://github.com/michaeliu4/wilf-s8/releases/tag/v1.0.0}{\texttt{github.com/michaeliu4/wilf-s8} (tag \texttt{v1.0.0})}.
The README distinguishes checks of retained evidence from full
regeneration. The companion contains the complete $\Sym_8$ certificate,
move provenance, the $9{,}510$ retained direct-placement recounts,
and the programs and inputs used for the structural checks.
% END RELEASE AVAILABILITY

\section*{AI statement}
The author and AI models made substantial mathematical contributions to this work. The author formulated the research programme around Questions 3.1, 3.7 and 3.8 of \cite{Vatter2026}, directed its stages, and set the verification requirements: independent recomputation of the classification by different algorithms, an explicit provenance for every class, and exhaustive machine checks of every structural statement. AI models of the Claude family (Anthropic; the versions are recorded in the author's project records) contributed, across the stages of the programme, the classification of $\Sym_8$ and its certificate, the interval formula and its derivation from the Ray--West classification, the catalogue of primitive returns, the substitution-tree deficit, the deletion-overlap formula, the layered formula and its proof, the occurrence complex and the universal constants, the design and implementation of the programs, and the computations. Of these, the cluster expansion, the deletion-graph analysis and \cref{thm:A}, the host-by-host identities of \cref{prop:weights,prop:chi2}, the Bessel-kernel formula of \cref{thm:bessel}, the independent recomputation of the classification, the programs of the companion archive, and the preparation and revision of this manuscript were contributed by Claude Fable~5.1 (19--21 September 2026). GPT-6 Astra Pro (OpenAI) contributed to
source verification, the subsequent manuscript revision, repairs to the
verification programs, the row-level recount, and the proofs of
\cref{prop:rectangle,cor:bessel-next}. The author takes full responsibility for all mathematical results and the contents of this manuscript.

\begin{thebibliography}{99}

\bibitem{BabsonWest2000}
E.~Babson and J.~West, \emph{The permutations $123p_4\cdots p_m$ and $321p_4\cdots p_m$ are Wilf-equivalent}, Graphs Combin. \textbf{16} (2000), no.~4, 373--380.

\bibitem{BWX2007}
J.~Backelin, J.~West and G.~Xin, \emph{Wilf-equivalence for singleton classes}, Adv. in Appl. Math. \textbf{38} (2007), no.~2, 133--148.

\bibitem{BHT2018}
D.~Bevan, C.~Homberger and B.~E. Tenner, \emph{Prolific permutations and permuted packings: downsets containing many large patterns}, J. Combin. Theory Ser. A \textbf{153} (2018), 98--121. Preprint: \href{https://arxiv.org/abs/1608.06931v3}{arXiv:1608.06931v3}.

\bibitem{BloomSaracino2014}
J.~Bloom and D.~Saracino, \emph{A simple bijection between $231$-avoiding and $312$-avoiding placements}, J. Combin. Math. Combin. Comput. \textbf{89} (2014), 23--32.

\bibitem{BHV2008}
R.~Brignall, S.~Huczynska and V.~Vatter, \emph{Simple permutations and algebraic generating functions}, J. Combin. Theory Ser. A \textbf{115} (2008), no.~3, 423--441.

\bibitem{BOO2000}
A.~Borodin, A.~Okounkov and G.~Olshanski, \emph{Asymptotics of Plancherel measures for symmetric groups}, J. Amer. Math. Soc. \textbf{13} (2000), no.~3, 481--515.

\bibitem{ESZ2015}
S.~B. Ekhad, N.~Shar and D.~Zeilberger, \emph{The number of $1\dots d$-avoiding permutations of length $d+r$ for SYMBOLIC $d$ but numeric $r$}, arXiv:1504.02513v1 (2015).

\bibitem{FRT1954}
J.~S. Frame, G.~de~B. Robinson and R.~M. Thrall, \emph{The hook graphs of the symmetric group}, Canad. J. Math. \textbf{6} (1954), 316--324.

\bibitem{Gessel1990}
I.~M. Gessel, \emph{Symmetric functions and P-recursiveness}, J. Combin. Theory Ser. A \textbf{53} (1990), no.~2, 257--285.

\bibitem{GouldenJackson1979}
I.~P. Goulden and D.~M. Jackson, \emph{An inversion theorem for cluster decompositions of sequences with distinguished subsequences}, J. London Math. Soc. (2) \textbf{20} (1979), no.~3, 567--576.

\bibitem{Hegarty2013}
P.~Hegarty, \emph{Permutations all of whose patterns of a given length are distinct}, J. Combin. Theory Ser. A \textbf{120} (2013), no.~7, 1663--1671, \url{https://doi.org/10.1016/j.jcta.2013.06.006}.

\bibitem{Homberger2012}
C.~Homberger, \emph{Counting fixed-length permutation patterns}, Online J. Anal. Comb. \textbf{7} (2012), Paper~3, 12~pp., doi:10.61091/ojac-703. Preprint: \href{https://arxiv.org/abs/1211.7117v1}{arXiv:1211.7117v1}.

\bibitem{Johansson2001}
K.~Johansson, \emph{Discrete orthogonal polynomial ensembles and the Plancherel measure}, Ann. of Math. (2) \textbf{153} (2001), no.~1, 259--296.

\bibitem{Kaplansky1945}
I.~Kaplansky, \emph{The asymptotic distribution of runs of consecutive elements}, Ann. Math. Statist. \textbf{16} (1945), no.~2, 200--203.

\bibitem{Krattenthaler2006}
C.~Krattenthaler, \emph{Growth diagrams, and increasing and decreasing chains in fillings of Ferrers shapes}, Adv. in Appl. Math. \textbf{37} (2006), no.~3, 404--431.

\bibitem{OEIS}
OEIS Foundation Inc., \emph{The On-Line Encyclopedia of Integer Sequences}, sequence A099952 (number of Wilf classes in $\Sym_n$), \url{https://oeis.org/A099952}.

\bibitem{Pratt1973}
V.~R. Pratt, \emph{Computing permutations with double-ended queues, parallel stacks and parallel queues}, in: Proceedings of the Fifth Annual ACM Symposium on Theory of Computing (STOC '73), ACM, 1973, pp.~268--277.

\bibitem{RayWest2003}
N.~Ray and J.~West, \emph{Posets of matrices and permutations with forbidden subsequences}, Ann. Comb. \textbf{7} (2003), no.~1, 55--88.

\bibitem{Schensted1961}
C.~Schensted, \emph{Longest increasing and decreasing subsequences}, Canad. J. Math. \textbf{13} (1961), 179--191.

\bibitem{Stankova1994}
Z.~E. Stankova, \emph{Forbidden subsequences}, Discrete Math. \textbf{132} (1994), no.~1--3, 291--316.

\bibitem{Stankova1996}
Z.~Stankova, \emph{Classification of forbidden subsequences of length 4}, European J. Combin. \textbf{17} (1996), no.~5, 501--517.

\bibitem{StankovaWest2002}
Z.~Stankova and J.~West, \emph{A new class of Wilf-equivalent permutations}, J. Algebraic Combin. \textbf{15} (2002), no.~3, 271--290.

\bibitem{Vatter2010}
V.~Vatter, \emph{Problems and conjectures}, in: Permutation Patterns (S.~Linton, N.~Ru\v{s}kuc and V.~Vatter, eds.), London Math. Soc. Lecture Note Ser., vol.~376, Cambridge Univ. Press, Cambridge, 2010, pp.~339--345, \url{https://doi.org/10.1017/CBO9780511902499.017}.

\bibitem{Vatter2026}
V.~Vatter, \emph{An assortment of problems in permutation patterns: unimodality, equivalence, derangements, and sorting}, arXiv:2602.16355v2 (2026).

\bibitem{Wolfowitz1944}
J.~Wolfowitz, \emph{Note on runs of consecutive elements}, Ann. Math. Statist. \textbf{15} (1944), no.~1, 97--98.

\end{thebibliography}

\end{document}
